\documentclass[12pt,a4paper,oneside]{report}
\usepackage[a4paper,left=1in,right=1in,top=1in,bottom=1in]{geometry}
% pdfLaTeX-compatible font setup (Latin Modern, close to Computer Modern)
\usepackage[T1]{fontenc}
\usepackage[utf8]{inputenc}
\usepackage{lmodern}
\usepackage{textcomp}
\usepackage{tfrupee}
\usepackage{amsmath,amssymb}
\usepackage{microtype}
\usepackage{graphicx}
\usepackage{float}
\usepackage{pdflscape}
\usepackage{caption}
\usepackage{booktabs}
% BlackBook grid-table formatting: uniform padding and single-weight rules throughout.
% Ordinary table header cells remain regular weight; only cost-section tables 4.6--4.10 use bold headers.
% Longtable data rows already end with \hline, so the last-foot rule is
% intentionally empty to prevent a doubled/heavier bottom border.
\renewcommand{\toprule}{\hline}
\renewcommand{\midrule}{\hline}
\renewcommand{\bottomrule}{}
\usepackage{tabularx}
\usepackage{longtable}
\usepackage{array}
\usepackage{calc}
\providecommand{\real}[1]{#1}
\usepackage{multirow}
\usepackage{setspace}
\usepackage{enumitem}
\usepackage{url}
\usepackage{xcolor}
\usepackage{newunicodechar}
\newunicodechar{≥}{\ensuremath{\geq}}
\newunicodechar{≠}{\ensuremath{\neq}}
% pdfLaTeX's utf8 input encoding already maps these characters, so they are
% only defined here for XeLaTeX/LuaLaTeX builds.
\usepackage{iftex}
\ifPDFTeX\else
  \newunicodechar{→}{\ensuremath{\rightarrow}}
  \newunicodechar{←}{\ensuremath{\leftarrow}}
  \newunicodechar{×}{\ensuremath{\times}}
  \newunicodechar{…}{\ldots}
  \newunicodechar{–}{--}
  \newunicodechar{—}{---}
  \newunicodechar{’}{'}
  \newunicodechar{“}{``}
  \newunicodechar{”}{''}
  \newunicodechar{«}{\guillemotleft}
  \newunicodechar{»}{\guillemotright}
\fi
% Rupee symbol is provided by tfrupee (command: \rupee).
\usepackage{listings}
\usepackage{chngcntr}
\usepackage{titlesec}
\usepackage{needspace}
\usepackage[hidelinks,unicode]{hyperref}
\usepackage{bookmark}
\counterwithin{figure}{chapter}
\counterwithin{table}{chapter}
\setlength{\parindent}{1.5em}
\setlength{\parskip}{0pt}
\setlength{\emergencystretch}{3em}
\renewcommand{\arraystretch}{1.18}
\setlength{\arrayrulewidth}{0.4pt}
\setlength{\extrarowheight}{2pt}
\setlength{\tabcolsep}{5pt}
\setcounter{secnumdepth}{3}
\setcounter{tocdepth}{3}
\titleformat{\chapter}[display]
  {\normalfont\bfseries\Huge}
  {\chaptername\ \thechapter}
  {0.6em}
  {\Huge}
\titlespacing*{\chapter}{0pt}{-0.2in}{1.6em}
\titleformat{\section}{\normalfont\bfseries\Large}{\thesection}{0.8em}{}
\titleformat{\subsection}{\normalfont\bfseries\large}{\thesubsection}{0.8em}{}
\titleformat{\subsubsection}{\normalfont\bfseries\normalsize}{\thesubsubsection}{0.8em}{}
\captionsetup[table]{position=top,skip=6pt}
% Equal space above and below every table caption:
\newcommand{\TableCapGap}{6pt}
\setlength{\LTpre}{0pt}
\setlength{\LTpost}{0pt}
\setlength{\intextsep}{\TableCapGap}
\captionsetup{font=small,labelfont=bf,justification=centering,hypcap=false,skip=6pt}
\pagestyle{plain}
\setlist[itemize]{leftmargin=2em,itemsep=2pt,topsep=3pt}
\setlist[enumerate]{leftmargin=2em,itemsep=2pt,topsep=3pt}
\lstdefinelanguage{Mermaid}{
  morekeywords={flowchart,sequenceDiagram,classDiagram,stateDiagram-v2,gantt,xychart-beta,pie,subgraph,end,participant,Note,alt,else,opt,loop,par,and,rect,autonumber},
  sensitive=true
}
\lstset{
  basicstyle=\ttfamily\footnotesize,
  columns=fullflexible,
  breaklines=true,
  breakatwhitespace=false,
  keepspaces=true,
  showstringspaces=false,
  tabsize=2,
  frame=single,
  rulecolor=\color{black!35},
  xleftmargin=0.5em,
  framexleftmargin=0.25em,
  aboveskip=0.7\baselineskip,
  belowskip=0.7\baselineskip,
  literate={→}{{$\rightarrow$}}1 {←}{{$\leftarrow$}}1 {–}{{--}}1 {—}{{---}}1 {’}{{'}}1 {“}{{``}}1 {”}{{''}}1 {…}{{\ldots}}1 {×}{{$\times$}}1 {«}{{\guillemotleft}}1 {»}{{\guillemotright}}1 {≠}{{$\neq$}}1 {≥}{{$\geq$}}1
}
\renewcommand{\arraystretch}{1.18}
\hypersetup{colorlinks=true,linkcolor=black,urlcolor=blue,citecolor=black}
\providecommand{\passthrough}[1]{#1}
\providecommand{\tightlist}{\setlength{\itemsep}{0pt}\setlength{\parskip}{0pt}}
% Display the supplied graphic when available; otherwise keep the report compilable
% and show a labeled placeholder until the original file is uploaded to Overleaf.
\newcommand{\SafeIncludeGraphics}[2][]{%
  \IfFileExists{#2}{\includegraphics[#1]{#2}}{%
    \par\noindent\begingroup\setlength{\fboxsep}{7pt}%
    \fbox{\begin{minipage}[c][0.15\textheight][c]{0.80\linewidth}
      \centering\textbf{Figure asset not found}\par\smallskip
      \texttt{\detokenize{#2}}\par\smallskip
      Upload this image file to the Overleaf project to display it.
    \end{minipage}}%
    \endgroup
  }%
}

% A wide diagram on its own portrait page, rotated 90 degrees so it fills the
% page. The page stays portrait; only the image and its caption are turned.
\newcommand{\RotatedFigure}[3]{%
  \clearpage
  \begin{figure}[p]
    \centering
    \rotatebox{90}{%
      \begin{minipage}{0.92\textheight}
        \centering
        \SafeIncludeGraphics[width=\linewidth,height=0.92\textwidth,keepaspectratio]{#1}%
        \par\medskip
        \captionof{figure}{#2}\label{#3}
      \end{minipage}%
    }%
  \end{figure}
  \clearpage
}

% A tall diagram on its own portrait page, filling the page height.
\newcommand{\TallFigure}[3]{%
  \clearpage
  \begin{figure}[p]
    \centering
    \SafeIncludeGraphics[width=\textwidth,height=0.9\textheight,keepaspectratio]{#1}%
    \par\medskip
    \captionof{figure}{#2}\label{#3}
  \end{figure}
  \clearpage
}

% Float placement: let large diagrams sit at the top of a text page and
% let text fill the remaining space instead of leaving blank pages.
\renewcommand{\topfraction}{0.9}
\renewcommand{\bottomfraction}{0.8}
\renewcommand{\textfraction}{0.07}
\renewcommand{\floatpagefraction}{0.8}
\setcounter{topnumber}{2}
\setcounter{bottomnumber}{2}
\setcounter{totalnumber}{4}

\begin{document}
\pagenumbering{roman}
\setcounter{page}{1}
\begin{titlepage}
\thispagestyle{empty}
\centering
\vspace*{0.55in}
{\normalsize A Project Report on\par}
\vspace{0.12in}
{\Large\bfseries VaultShield: Sandboxed, Monitored Execution\\\par}
{\Large\bfseries and Human Review for AI Coding Agents\par}
\vspace{0.38in}
{\normalsize Submitted in partial fulfillment of the requirements for the award\\ of the degree of\par}
\vspace{0.34in}
{\large\bfseries Bachelor of Engineering\par}
\vspace{0.12in}
{\normalsize in\par}
\vspace{0.12in}
{\large\bfseries Computer Science and Engineering (Data Science)\par}
\vspace{0.30in}
{\normalsize by\par}
\vspace{0.12in}
{\large\bfseries
Aayush Subhash Balip (24207010)\\
Krish Mahesh Patil (24207006)\\
Muhammad Anis Momin (24207008)\\
Sameer Hasan Saiyed (24207014)\par}
\vspace{0.36in}
{\normalsize Under the Guidance of\par}
\vspace{0.08in}
{\large\bfseries Prof. Rajashri Chaudhari\\ Prof. Hardiki Achrekar\par}
\vfill
% Institute logo: upload it to the Overleaf project as img-000.jpg.
\IfFileExists{img-000.jpg}{\SafeIncludeGraphics[width=1.65in]{img-000.jpg}\par}{}
\vspace{0.16in}
{\large\bfseries Computer Science and Engineering (Data Science)\par}
{\normalsize A.P. Shah Institute of Technology\par}
{\normalsize G.B. Road, Kasarvadavali, Thane (W) -- 400615\par}
{\normalsize UNIVERSITY OF MUMBAI\par}
\vspace{0.08in}
{\normalsize\bfseries Academic Year 2026--2027\par}
\end{titlepage}

% ------------------------------------------------------------------
% APPROVAL SHEET — layout follows the NirmalSagar BlackBook template.
% Keep the approval paragraph justified/left aligned; signature blocks
% are two unboxed columns; the HOD appears centered below them.
% ------------------------------------------------------------------
\thispagestyle{empty}
\vspace*{0.55in}
\begin{center}\large\bfseries Approval Sheet\end{center}
\vspace{0.35in}
\noindent This Project Report entitled \textbf{\textit{VaultShield: Sandboxed, Monitored Execution and Human Review for AI Coding Agents}} submitted by \textbf{Aayush Subhash Balip} (24207010), \textbf{Krish Mahesh Patil} (24207006), \textbf{Muhammad Anis Momin} (24207008), and \textbf{Sameer Hasan Saiyed} (24207014) is approved for the partial fulfillment of the requirement for the award of the degree of \textbf{Bachelor of Engineering in Computer Science and Engineering (Data Science)} from \textbf{University of Mumbai}.\par

\vspace{0.95in}
\noindent\begin{tabularx}{\textwidth}{>{\centering\arraybackslash}X >{\centering\arraybackslash}X}
Prof. Rajashri Chaudhari & Prof. Hardiki Achrekar\\[0.08in]
Guide & Co-Guide
\end{tabularx}

\vspace{1.15in}
\begin{center}
Dr. Pravin Adivarekar\\[0.08in]
HOD, Computer Science and Engineering (Data Science)
\end{center}
\vfill
\noindent Place: A.P. Shah Institute of Technology, Thane\\[0.15in]
Date:
\clearpage

% ------------------------------------------------------------------
% CERTIFICATE — two unboxed signature columns, HOD/Principal below,
% and unboxed examiner entries as in the NirmalSagar template.
% ------------------------------------------------------------------
\thispagestyle{empty}
\vspace*{0.55in}
\begin{center}\large\bfseries CERTIFICATE\end{center}
\vspace{0.35in}
\noindent This is to certify that the project entitled \textbf{\textit{VaultShield: Sandboxed, Monitored Execution and Human Review for AI Coding Agents}} submitted by \textbf{Aayush Subhash Balip} (24207010), \textbf{Krish Mahesh Patil} (24207006), \textbf{Muhammad Anis Momin} (24207008), and \textbf{Sameer Hasan Saiyed} (24207014) for the partial fulfillment of the requirement for award of a degree of \textbf{Bachelor of Engineering in Computer Science and Engineering (Data Science)} to the University of Mumbai, is a bonafide work carried out during academic year 2026--2027.\par

\vspace{0.78in}
\noindent\begin{tabularx}{\textwidth}{>{\centering\arraybackslash}X >{\centering\arraybackslash}X}
\textbf{Prof. Rajashri Chaudhari} & \textbf{Prof. Hardiki Achrekar}\\[0.08in]
Guide & Co-Guide
\end{tabularx}

\vspace{0.95in}
\noindent\begin{tabularx}{\textwidth}{>{\centering\arraybackslash}X >{\centering\arraybackslash}X}
\textbf{Dr. Pravin Adivarekar} & \textbf{Dr. Uttam D. Kolekar}\\[0.08in]
HOD, C.S.E (Data Science) & Principal
\end{tabularx}

\vspace{0.72in}
\noindent\begin{tabularx}{\textwidth}{>{\raggedright\arraybackslash}X >{\raggedright\arraybackslash}X}
External Examiner(s) & Internal Examiner(s)\\[0.20in]
1. & 1.\\[0.34in]
2. & 2.
\end{tabularx}

\vfill
\noindent Place: A.P. Shah Institute of Technology, Thane\\[0.15in]
Date:
\clearpage

\thispagestyle{empty}
\vspace*{0.55in}
\begin{center}\large\bfseries Acknowledgement\end{center}
\vspace{0.55in}
We express our sincere gratitude to our project guide, \textbf{Prof. Rajashri Chaudhari}, for the constant guidance, encouragement and critical feedback provided throughout this project. The guide's insistence on precision, both in engineering and in stating the limits of what had been built, shaped the security model at the heart of this work.

We are equally grateful to our co-guide, \textbf{Prof. Hardiki Achrekar}, for valuable suggestions on system design and evaluation, and for the time spent reviewing our progress at every stage.

We thank \textbf{Dr. Pravin Adivarekar}, Head of the Department of Computer Science and Engineering (Data Science), for providing the facilities and the supportive environment that made this project possible, and \textbf{Dr. Uttam D. Kolekar}, Principal, A. P. Shah Institute of Technology, for the institutional support extended to us.

We also thank the project coordinators for organising the reviews and keeping us on schedule, and all the teaching and non-teaching staff of the department for their help throughout the academic year.

Finally, we thank our colleagues and friends for their feedback on the dashboard and the documentation, and our families for their patience and encouragement.

\vspace{0.45in}
\begin{flushleft}
\textbf{Aayush Subhash Balip}\\(24207010)\\[0.15in]
\textbf{Krish Mahesh Patil}\\(24207006)\\[0.15in]
\textbf{Muhammad Anis Momin}\\(24207008)\\[0.15in]
\textbf{Sameer Hasan Saiyed}\\(24207014)
\end{flushleft}
\clearpage

\thispagestyle{empty}
\vspace*{0.55in}
\begin{center}\large\bfseries Declaration\end{center}
\vspace{0.35in}
We declare that this written submission represents our ideas in our own words, and where others' ideas or words have been included, we have adequately cited and referenced the original sources. We also declare that we have adhered to all principles of academic honesty and integrity and have not misrepresented, fabricated or falsified any idea, data, fact or source in our submission. We understand that any violation of the above will be cause for disciplinary action by the Institute and can also evoke penal action from the sources which have thus not been properly cited or from whom proper permission has not been taken when needed.

\vspace{0.5in}
\begin{flushright}
(Signature)\\ Aayush Subhash Balip (24207010)\\[0.28in]
(Signature)\\ Krish Mahesh Patil (24207006)\\[0.28in]
(Signature)\\ Muhammad Anis Momin (24207008)\\[0.28in]
(Signature)\\ Sameer Hasan Saiyed (24207014)
\end{flushright}
\vfill
Date: 
\clearpage

\thispagestyle{empty}
\vspace*{1.60in}
\begin{center}{\large\bfseries Abstract}\end{center}
\vspace{0.20in}
AI coding agents now read repositories, edit many files and execute shell commands with little human intervention. Run directly on a developer's machine, such an agent inherits the developer's privileges: it can read real secrets, write anywhere, contact any host and leave unreviewed changes in version control. Current tooling conflates running an agent with trusting its output. This project presents VaultShield, a pipeline that separates the two. VaultShield runs the agent in a disposable Docker container holding a fresh clone of the repository and a read-only credential file whose values mirror the shape of the developer's \texttt{.env} but contain no real secret; a repository-host token, when needed, is confined to the single \texttt{git clone} process. Terminal output, in-container file changes and policy violations are merged into a timestamped JSON Lines event log. Writes to off-limits paths fail the run, and an optional egress policy routes traffic through a hardened proxy that blocks and records non-allowlisted hosts. A separate hardened helper container then rebuilds the repository from a pre-agent base and commits the change to a dedicated branch with a structured diff, so the agent cannot falsify the evidence. Artifacts are uploaded to a Django REST Framework backend on MongoDB, scanned asynchronously, and approved or rejected by a reviewer in a React dashboard. Only an approved run can become a pull or merge request. The system comprises about 18,100 logical lines of product code. Of 2,704 automated tests, including Docker-backed isolation, egress, token-hygiene and tamper-resistance tests and concurrency tests against a real MongoDB replica set, 2,702 passed. In a 40-run benchmark, VaultShield added 1.17 seconds of fixed overhead per run by default and 1.98 seconds with egress allowlisting, a projected 0.4--0.7\% of a five-minute agent session; the host checkout was untouched and the change was captured in every run, with no container or network left behind. The evaluation also exposed limits: the shared-kernel boundary of containers, opt-in egress control, a polling file watcher that missed 20\% of very short writes, and an approval gate that is a workflow control rather than a separation-of-duties control in the single-user MVP. VaultShield shows that containment, evidence and human review for coding agents can be delivered together at negligible cost, with every residual risk stated explicitly.

\vspace{0.15in}
\noindent\textbf{Keywords:} AI coding agents, sandboxing, Docker, human-in-the-loop review, audit logging, credential isolation, egress control, software supply-chain security.
\clearpage

\tableofcontents
\clearpage
\listoffigures
\clearpage
\listoftables
\clearpage

\chapter*{List of Abbreviations}\label{list-of-abbreviations}
\begin{longtable}{@{}p{1.1in}@{}p{0.25in}@{}p{4.9in}@{}}
AA &: & Level AA conformance of the Web Content Accessibility Guidelines \\[2pt]
ACI &: & Agent-Computer Interface \\[2pt]
API &: & Application Programming Interface \\[2pt]
ASGI &: & Asynchronous Server Gateway Interface \\[2pt]
CI/CD &: & Continuous Integration / Continuous Delivery \\[2pt]
CLI &: & Command-Line Interface \\[2pt]
COCOMO &: & Constructive Cost Model \\[2pt]
CORS &: & Cross-Origin Resource Sharing \\[2pt]
CPU &: & Central Processing Unit \\[2pt]
CWE &: & Common Weakness Enumeration \\[2pt]
DNS &: & Domain Name System \\[2pt]
DRF &: & Django REST Framework \\[2pt]
HTTP &: & Hypertext Transfer Protocol \\[2pt]
HTTPS &: & Hypertext Transfer Protocol Secure \\[2pt]
IDE &: & Integrated Development Environment \\[2pt]
IP &: & Internet Protocol \\[2pt]
JSON &: & JavaScript Object Notation \\[2pt]
JSONL &: & JSON Lines \\[2pt]
KLOC &: & Thousand Lines of Code \\[2pt]
LLM &: & Large Language Model \\[2pt]
LM &: & Language Model \\[2pt]
MR &: & Merge Request \\[2pt]
MVP &: & Minimum Viable Product \\[2pt]
OAuth &: & Open Authorization \\[2pt]
OCR &: & Optical Character Recognition \\[2pt]
PM &: & Person-Month \\[2pt]
PR &: & Pull Request \\[2pt]
PTY &: & Pseudo-Terminal \\[2pt]
RBAC &: & Role-Based Access Control \\[2pt]
REST &: & Representational State Transfer \\[2pt]
SHA &: & Secure Hash Algorithm \\[2pt]
SLOC &: & Source Lines of Code \\[2pt]
SNI &: & Server Name Indication \\[2pt]
SPA &: & Single-Page Application \\[2pt]
SQL &: & Structured Query Language \\[2pt]
STRIDE &: & Spoofing, Tampering, Repudiation, Information disclosure, Denial of service, Elevation of privilege \\[2pt]
TLS &: & Transport Layer Security \\[2pt]
TTL &: & Time To Live \\[2pt]
UI &: & User Interface \\[2pt]
UML &: & Unified Modeling Language \\[2pt]
URL &: & Uniform Resource Locator \\[2pt]
UTC &: & Coordinated Universal Time \\[2pt]
VCS &: & Version Control System \\[2pt]
VM &: & Virtual Machine \\[2pt]
WCAG &: & Web Content Accessibility Guidelines \\[2pt]
WSGI &: & Web Server Gateway Interface \\[2pt]
\end{longtable}

\clearpage
\pagenumbering{arabic}
\setcounter{page}{1}

\chapter{Introduction}\label{introduction}

Software development is in the middle of a structural change. Large Language Models trained on source code have moved from autocompleting a single line to operating as autonomous coding agents: programs that read a repository, plan a change, edit many files, execute shell commands, run test suites and iterate on the results with little or no human intervention between steps. Tools such as OpenAI Codex CLI, Google Antigravity, Amazon's Kiro and Anthropic's Claude Code expose this capability directly in the developer's terminal or editor. Research systems such as SWE-agent \cite{ref3} and OpenHands \cite{ref4} have demonstrated that such agents can resolve a meaningful fraction of real GitHub issues drawn from benchmarks like SWE-bench \cite{ref2}. The productivity case is strong, and adoption in industry and academia is accelerating accordingly.

The same capability that makes coding agents useful also makes them hazardous. Operationally, an agent is an ordinary operating-system process that runs with the privileges of the user who launched it. When it is started directly on a developer's workstation, it can read every secret in the shell environment and in nearby .env files, write anywhere the user can write, execute commands the prompt never asked for, make outbound network connections to arbitrary hosts, and leave changes in the version control system before anyone has examined them. Prior work has shown that LLM-generated code is frequently insecure \cite{ref8}, that LLM-integrated applications can be hijacked through indirect prompt injection embedded in the data they process \cite{ref7}, and that even the most carefully aligned agents take risky actions in a non-trivial fraction of high-stakes scenarios \cite{ref5}. An agent that reads a malicious README or a poisoned dependency is therefore not a hypothetical concern; it is a direct consequence of giving a probabilistic system unsupervised access to a real machine.

Most agent tooling today conflates two distinct questions: running the agent and trusting its output. Once the agent has been started on the developer's machine, its file writes, its network calls and its use of credentials have already happened, and the only remaining safeguard is whatever review the developer chooses to perform afterwards on a working tree that the agent has already modified. Governance guidance for agentic AI systems explicitly recommends constraining an agent's action space, requiring human approval for consequential actions, making agent activity legible, and keeping systems interruptible \cite{ref9}. Related work argues that meaningful oversight of agents depends on visibility: identifiers, real-time monitoring and activity logs \cite{ref10}. These recommendations are widely cited, but they are rarely packaged as a concrete, reproducible pipeline that a developer can run against an existing repository.

This project, VaultShield, addresses that gap. VaultShield runs an AI coding agent inside a disposable Docker container, records what the agent does, captures the code changes it produces as a reviewable Git diff on a dedicated branch, and holds those changes behind an explicit human approval step before they can become a pull request or merge request. The developer's working tree, real credentials and repository write access are never handed to the agent. The system consists of a Python command-line interface that orchestrates the sandboxed run on the developer's host, a per-run Docker execution plane, and a Django REST Framework backend with a React dashboard where a reviewer inspects evidence and approves or rejects the run. The result is a pipeline that separates execution from trust: producing evidence never implies accepting it.

The remainder of this chapter presents the motivation for the work, the formal problem statement, the measurable objectives the project set out to achieve, and the scope and explicit exclusions of the delivered system. Chapter 2 reviews the literature on which the design draws, and Chapters 3 through 7 describe the design, planning, implementation, testing and evaluation of the system.

\section{Motivation}\label{motivation}

The first motivation is the asymmetry between the cost of an agent mistake and the cost of preventing it. An agent that deletes the wrong directory, overwrites a configuration file or commits a hard-coded key costs the developer hours of recovery, and in the case of leaked credentials may cost an organisation far more. By contrast, running the same agent in a throwaway container costs, as measured in Chapter 7, roughly 1.2 seconds of fixed overhead per run on the evaluation machine. When the protective step is this cheap relative to the potential damage, the absence of such a step in everyday tooling is difficult to justify, and a practical, ready-to-use pipeline becomes valuable.

The second motivation is credential exposure. Modern applications cannot start without configuration: database connection strings, API keys and session secrets are typically stored in a local .env file that sits in the very repository the agent is asked to modify. An agent working directly on the host can read that file, echo its contents into a terminal transcript, or embed a value in generated code. A preliminary spike conducted during this project confirmed this directly: with no protection in place, a production coding agent asked to reveal a value from .env read the file with its built-in file tool and printed both test secrets in its answer. Agent-specific deny rules reduced but did not eliminate the exposure, and only removing the real value from disk proved reliable. VaultShield therefore treats secret handling as a first-class design concern rather than an afterthought.

The third motivation is reviewability and accountability. Code review is the primary quality-control mechanism of modern software engineering \cite{ref14}, and developers already understand the pull-request workflow: a change is proposed on a branch, a diff is examined, and a human decides whether to merge. Agent output, however, frequently bypasses this workflow because it lands directly in the developer's working tree. Capturing agent output on a dedicated branch with a structured diff, accompanied by a timestamped event log and screenshots of the running application, lets a reviewer apply the same mental model to agent changes that they already apply to human changes. This also produces a durable audit trail that answers the question ``what did the agent actually do?'' after the fact.

The fourth motivation is honest security engineering. Many tools describe themselves as ``secure sandboxes'' without stating which controls they implement and which risks remain. Container isolation, in particular, shares the host kernel and is not equivalent to a virtual machine boundary \cite{ref11}, \cite{ref12}, \cite{ref13}. A further aim of this project was therefore to build a system whose security model is explicit: every control is traceable to a specific module, and every residual risk is documented rather than hidden. This honesty is itself a contribution, because it allows adopters to make informed decisions about where VaultShield fits in their own threat model.

Finally, the project was motivated by its educational value for a final-year engineering team. Building VaultShield required working across operating-system isolation, version control internals, distributed backend design, front-end engineering and empirical evaluation. This breadth aligns closely with the learning outcomes of a Computer Science and Engineering programme and provided substantial exposure to production-grade engineering practices such as property-based testing and threat modelling.

\section{Problem Statement}\label{problem-statement}

Given a Git repository, a natural-language prompt and an AI coding agent, the problem is to allow the agent to attempt the requested change while guaranteeing that the agent never executes against the developer's real working tree, real application secrets or repository-host credentials, everything the agent does is recorded as reviewable evidence, the change it produces is captured faithfully in a form the agent could not have tampered with, and no change reaches a shared branch or pull request without an explicit human approval decision.

Stated as an engineering problem, the system must accept a repository URL, a prompt and an agent identifier; create an isolated execution environment containing a fresh clone of the repository and a dummy credential file; run the agent under monitoring; capture terminal output, file changes, network egress decisions and screenshots; build a branch and structured diff from a trusted base; upload the artifacts to a backend; and expose them through a dashboard on which a reviewer approves or rejects the run. Only an approved run may be submitted as a PR or MR, and every container, network and temporary branch must be removed on every exit path, including failures, timeouts and user interruption.

The problem is constrained by practical requirements. The agent must still be able to work: it needs a real repository, a functioning toolchain and, in some cases, its own model API credential. The application under change must be able to start inside the sandbox for screenshots, which means that configuration values must keep their structural shape even though their content is fake. The pipeline overhead must remain small compared with typical agent run times so that developers are not discouraged from using it. Finally, the system must be candid about the limits of container isolation, because overstating its guarantees would itself create risk.

\section{Objectives}\label{objectives}

The project set out to achieve the following objectives.
\begin{enumerate}
\def\labelenumi{\arabic{enumi}.}
\tightlist
\item
  Develop a universal AI agent execution platform compatible with multiple AI coding agents and development environments through a common CLI interface.
\item
  Implement an intelligent Agentic AI Orchestrator that understands project context, asks clarification questions before execution, and provides enriched context to AI agents to improve accuracy and reduce hallucinations.
\item
  Ensure secure AI-assisted software development by executing AI agents inside isolated Docker-based sandbox environments while protecting production credentials and resources.
\item
  Provide comprehensive monitoring and automated security analysis by capturing AI activities, execution logs, code changes, and identifying security vulnerabilities or unsafe operations.
\item
  Enable transparent governance and human approval through a web-based dashboard that allows developers to review, verify, approve, or reject AI-generated changes before deployment.
\end{enumerate}

\section{Scope}\label{scope}

The scope of the VaultShield project covers the sandboxed execution, monitoring and human review of AI coding agents working on a Git repository. The project focuses on keeping the agent away from the developer's real working tree and secrets while producing evidence that a reviewer can examine before any change is accepted.

\begin{itemize}
\item
  \textbf{Command-Line Interface:}\\
  The project provides a host-side vaultshield CLI with the init, login, run, submit and teardown commands, together with the supporting ask and protect command groups.
\item
  \textbf{Sandboxed Execution:}\\
  Each run uses a per-run Docker sandbox with a fresh clone of the repository, a read-only dummy credential file and guaranteed teardown.
\item
  \textbf{Agent Adapters:}\\
  The mock, codex, kiro and antigravity agents are supported behind a common adapter interface.
\item
  \textbf{Monitoring and Logging:}\\
  The system captures terminal output, watches file changes inside the container, enforces off-limits paths and records a JSON Lines event log.
\item
  \textbf{Network Egress Control:}\\
  An optional per-run egress policy is enforced by a hardened proxy sidecar on an internal Docker network.
\item
  \textbf{Trusted Change Capture:}\\
  The run branch and structured diff are built in a separate hardened helper container, and run outputs are scanned for leaked credentials.
\item
  \textbf{Backend and Review State Machine:}\\
  A Django REST Framework backend on MongoDB provides token authentication, the run state machine, artifact ingest, and Celery-based security scanning and cleanup.
\item
  \textbf{Review Dashboard:}\\
  A React dashboard lets a reviewer browse projects and runs, inspect the evidence and record approve or reject decisions.
\item
  \textbf{Pull and Merge Request Submission:}\\
  Approved runs can be submitted as pull requests on GitHub or merge requests on GitLab.
\end{itemize}

The following are outside the scope of the project:

\begin{itemize}
\item
  \textbf{Hosted Execution:}\\
  Running the CLI inside a container or as a hosted service is not covered.
\item
  \textbf{Virtual-Machine Isolation:}\\
  Virtual-machine-grade isolation is not provided; container escape is a documented residual risk.
\item
  \textbf{Non-CLI Agents:}\\
  Agents that cannot be invoked as a command-line process inside a Linux container are not supported.
\item
  \textbf{Intent Analysis:}\\
  The system records agent behaviour but does not judge agent intent automatically.
\item
  \textbf{Content-Level Filtering:}\\
  Network filtering is host-based; URL-path-level or content-level filtering is not performed.
\item
  \textbf{Encoded Exfiltration:}\\
  Deliberately encoded exfiltration and credentials rendered inside screenshots are not detected.
\item
  \textbf{Multi-User Operation:}\\
  Multi-tenant operation, role-based access control and separate runner and reviewer identities are not included.
\item
  \textbf{Live Streaming:}\\
  Agent output is not streamed live during a run.
\item
  \textbf{Other Integrations:}\\
  Other repository hosts, automatic merging and CI/CD pipeline integration are not included.
\end{itemize}


The minimum viable product is scoped to a single developer's machine. The backend data model intentionally contains no owner, team or tenant fields, and the dashboard's sign-up page is a placeholder; accounts are created through the Django administration interface. The approval gate is therefore a workflow control within one developer's trust domain rather than a security boundary between mutually distrustful parties; this report discusses the consequences of this decision later. The project also does not claim to be a formally verified security boundary. Its contribution is to reduce the blast radius of AI-assisted development and to make agent activity observable and reviewable, with every remaining risk stated explicitly.

\clearpage

\chapter{Literature Review}\label{literature-review}

This chapter surveys the research on which VaultShield builds. The literature falls into five overlapping strands: the capabilities and evaluation of autonomous coding agents, the security risks of LLM-generated code and LLM agents, isolation technologies for running untrusted code, monitoring and governance of agentic systems, and human code review. Each strand is reviewed below, followed by a comparative analysis and a statement of the research gap that motivates the project. Citations follow the IEEE numbered style, and the full references appear in the Bibliography.

Code-generating models and the need for execution sandboxes. Chen et al.~\cite{ref1} introduced Codex, a GPT model fine-tuned on public GitHub code, together with the HumanEval benchmark of hand-written programming problems evaluated by functional correctness. Because evaluating functional correctness requires executing model-generated programs, the authors describe running those programs inside a sandboxed environment based on gVisor, and they discuss misuse, over-reliance and the generation of insecure code as broader risks. The paper is significant for VaultShield because it established, at the very beginning of the code-LLM era, that model output must be treated as untrusted code to be executed only inside an isolation boundary. Its sandbox, however, was an internal evaluation harness for short, self-contained functions, not a developer-facing tool for repository-scale changes with human review.

Repository-scale agents and their evaluation. Jimenez et al.~\cite{ref2} proposed SWE-bench, a benchmark of 2,294 task instances drawn from real GitHub issues and pull requests across twelve popular Python repositories, where a model must produce a patch that makes the repository's failing tests pass. SWE-bench moved evaluation from isolated functions to whole repositories and showed that the best model at the time resolved fewer than 2\% of issues, demonstrating how much harder realistic software engineering is. Yang et al.~\cite{ref3} then introduced SWE-agent, arguing that agents need an agent-computer interface with purpose-built commands for viewing, searching and editing files; with GPT-4 Turbo it resolved 12.5\% of SWE-bench tasks, and it executes inside a Docker container. Wang et al.~\cite{ref4} presented OpenHands, an open platform in which agents act through a sandboxed Docker runtime and every action and observation is recorded in an event stream. These systems validate two of VaultShield's core design choices: that a container is the practical execution environment for repository-scale agents, and that a chronological record of agent actions is the natural unit of observability. Their focus, however, is on agent capability: the sandbox exists to give the agent a reproducible workspace, and none of them provide a gate that holds the agent's patch for human approval or a credential model that keeps the developer's real secrets out of the workspace.

Risks of LLM agents. Ruan et al.~\cite{ref5} introduced ToolEmu, which uses an LM to emulate tool execution so that agents can be red-teamed against high-stakes scenarios without real side effects, together with an automatic safety evaluator. Across 36 toolkits and 144 test cases, they found that even the safest agent they evaluated exhibited risky failures in 23.9\% of cases. Debenedetti et al.~\cite{ref6} built AgentDojo, a dynamic environment with 97 realistic tasks and 629 security test cases for measuring prompt-injection attacks against tool-using agents and the defences against them, and showed that both attacks and defences remain far from solved. Greshake et al.~\cite{ref7} demonstrated indirect prompt injection: instructions planted in data that an LLM-integrated application retrieves can take control of the application and cause data exfiltration or other harmful actions. For coding agents the implication is direct, because a repository's own files, issues and dependencies are untrusted input. These works characterise the threat convincingly, but they are evaluation frameworks; they do not provide an operational containment and review pipeline for developers.

Security of AI-generated code. Pearce et al.~\cite{ref8} systematically prompted GitHub Copilot with 89 scenarios relevant to MITRE's top Common Weakness Enumeration categories and found that approximately 40\% of the 1,689 generated programs contained vulnerabilities. This result motivates VaultShield's static security scan over every ingested diff and, more fundamentally, its insistence that a human reviews agent code before it is merged. The study evaluated completion suggestions rather than autonomous agents, so it does not address agents that also execute commands; VaultShield extends the concern from ``insecure code'' to ``insecure behaviour''.

Governance and visibility of agentic systems. Shavit et al.~\cite{ref9} proposed a set of practices for keeping agentic AI systems safe and accountable, including evaluating suitability for a task, constraining the action space and requiring human approval for high-impact actions, setting safe defaults, making agent activity legible to users, automatic monitoring, attributability and interruptibility. Chan et al.~\cite{ref10} argued that oversight of AI agents requires visibility and identified three categories of measures: agent identifiers, real-time monitoring and activity logging. VaultShield can be read as a concrete implementation of several of these practices for the coding-agent setting: the run identifier is an agent-activity identifier that names the sandbox, the branch, the run directory and the backend record; the event log, egress log and screenshots constitute activity logging; and the approval state machine is the human-approval control. Both papers are deliberately technology-neutral policy proposals, so they leave the engineering of such controls as an open problem.

Container isolation and its limits. Combe et al.~\cite{ref11} analysed Docker from a security perspective and identified the main areas of concern: the isolation provided by Linux namespaces and control groups, the shared host kernel, network configuration, and the trustworthiness of the image distribution chain. They concluded that Docker's default configuration offers a reasonable level of isolation but that its security depends heavily on configuration and on the ecosystem around it. Agache et al.~\cite{ref12} described Firecracker, the virtual machine monitor behind AWS Lambda, which provides hardware-virtualisation isolation with microVMs that boot in as little as 125 ms with less than 5 MiB of memory overhead, demonstrating that VM-grade isolation can approach container-like efficiency. Young et al.~\cite{ref13} measured gVisor, a user-space kernel that intercepts system calls to reduce the host kernel's attack surface, and showed that this stronger isolation carries substantial performance costs for system-call-heavy and file-system-heavy workloads. Together these papers frame VaultShield's isolation trade-off: Docker containers are cheap and universally available but share the kernel, while gVisor and microVMs offer stronger boundaries at a cost in performance or operational complexity. VaultShield chooses Docker for its MVP and documents the escape risk explicitly, identifying gVisor and microVMs as future work.

Human code review. Bacchelli and Bird \cite{ref14} studied modern code review at Microsoft through observation, interviews and surveys. They found that although finding defects is the primary stated motivation for review, much of its actual value lies in knowledge transfer, team awareness and alternative solutions, and that the dominant challenge for reviewers is understanding the change and its context. This finding directly shapes VaultShield's dashboard: a reviewer is given not only the diff but also the timeline of what the agent did, screenshots of the running application and automated findings, so that understanding an agent's change does not require reconstructing its reasoning from the patch alone.

\section{Comparative Analysis of Recent Studies}\label{comparative-analysis-of-recent-studies}

Table 2.1 compares the reviewed works along the dimensions most relevant to this project and records, for each, the limitation that VaultShield addresses.

\begingroup\small

\Needspace{8\baselineskip}
\par\vspace{\TableCapGap}
\noindent\begin{minipage}{\textwidth}
\captionof{table}{Comparative analysis of the reviewed literature}\label{tab:2-1}
\end{minipage}
\par
\begin{longtable}[]{@{}|
  >{\raggedright\arraybackslash}p{(\columnwidth - 11\tabcolsep) * \real{0.2000}} |
  >{\raggedright\arraybackslash}p{(\columnwidth - 11\tabcolsep) * \real{0.2000}} |
  >{\raggedright\arraybackslash}p{(\columnwidth - 11\tabcolsep) * \real{0.2000}} |
  >{\raggedright\arraybackslash}p{(\columnwidth - 11\tabcolsep) * \real{0.2000}} |
  >{\raggedright\arraybackslash}p{(\columnwidth - 11\tabcolsep) * \real{0.2000}}|@{}}
\hline
Ref. & Approach & Key techniques & Tools / platform & Limitation relative to this project \\
\hline
\endfirsthead
\hline
Ref. & Approach & Key techniques & Tools / platform & Limitation relative to this project \\
\hline
\endhead
\endlastfoot
\cite{ref1} Chen et al., 2021 & Code LLM with functional-correctness evaluation & Fine-tuning on GitHub code; pass@k; sandboxed execution of generated programs & Codex, HumanEval, gVisor-based sandbox & Sandbox is an internal evaluation harness for short functions; no repository workflow, credential model or human review \\\hline
\cite{ref2} Jimenez et al., 2024 & Repository-level benchmark from real issues & Issue-to-patch tasks validated by repository test suites & SWE-bench & Measures capability only; patches are applied automatically for scoring, not held for approval \\\hline
\cite{ref3} Yang et al., 2024 & Agent with a purpose-built agent-computer interface & Custom file viewing, search and edit commands; containerised execution & SWE-agent, Docker, GPT-4 Turbo & Container is a workspace, not a trust boundary; no secret substitution, egress control or review gate \\\hline
\cite{ref4} Wang et al., 2025 & Open platform for generalist software agents & Event-stream architecture; sandboxed runtime; multiple agent implementations & OpenHands, Docker runtime & Focus on agent capability and extensibility; changes are not captured on a reviewable branch behind an approval state machine \\\hline
\cite{ref5} Ruan et al., 2024 & LM-emulated sandbox for red-teaming agents & Emulated tool execution; automatic safety evaluator & ToolEmu & Emulation avoids real side effects but cannot contain a real agent acting on a real repository \\\hline
\cite{ref6} Debenedetti et al., 2024 & Dynamic benchmark for prompt-injection attacks and defences & Tool-using agent tasks with injected adversarial content & AgentDojo & Evaluates attacks and defences; offers no operational containment or audit pipeline \\\hline
\cite{ref7} Greshake et al., 2023 & Attack study on LLM-integrated applications & Indirect prompt injection via retrieved data & Real LLM-integrated applications & Characterises the threat; mitigation is left open, motivating containment and review \\\hline
\cite{ref8} Pearce et al., 2022 & Empirical security study of AI code suggestions & CWE-guided scenarios; static analysis of generated code & GitHub Copilot, CodeQL & Studies code suggestions, not autonomous agents that also execute commands \\\hline
\cite{ref9} Shavit et al., 2023 & Governance practices for agentic AI & Approval for high-impact actions; legibility; monitoring; interruptibility & Policy white paper & Technology-neutral; no reference implementation for coding agents \\\hline
\cite{ref10} Chan et al., 2024 & Visibility measures for AI agents & Agent identifiers; real-time monitoring; activity logging & Policy and analysis paper & Defines what to log, not how to capture tamper-resistant evidence from a running agent \\\hline
\cite{ref11} Combe et al., 2016 & Security analysis of Docker & Namespaces, cgroups, capabilities, network and image-chain review & Docker & Shows security depends on configuration; does not address agent-specific credential or review concerns \\\hline
\cite{ref12} Agache et al., 2020 & Lightweight VM isolation for multi-tenant serverless & MicroVMs with minimal device model & Firecracker, KVM & Stronger isolation but requires KVM and different tooling; not integrated with a developer review workflow \\\hline
\cite{ref13} Young et al., 2019 & Performance study of a user-space kernel sandbox & System-call interception; overhead measurement & gVisor & Quantifies the cost of stronger isolation; informs, but is not, an agent pipeline \\\hline
\cite{ref14} Bacchelli \& Bird, 2013 & Empirical study of modern code review & Observation, interviews, surveys & Microsoft CodeFlow & Concerns human-authored changes; review of agent changes needs additional behavioural context \\\hline\end{longtable}
\addtocounter{table}{-1}
\par\medskip

\endgroup

\section{Research Gap}\label{research-gap}

The reviewed literature establishes four facts with considerable consistency. First, coding agents are now capable enough to make repository-scale changes \cite{ref2}, \cite{ref3}, \cite{ref4}, so they will increasingly be given real repositories. Second, both the code they generate and the actions they take are risky in measurable ways \cite{ref5}, \cite{ref6}, \cite{ref7}, \cite{ref8}. Third, containers are the de facto execution environment for such agents, but container security depends on configuration and does not equal a VM boundary \cite{ref11}, \cite{ref12}, \cite{ref13}. Fourth, governance research recommends human approval, legible activity logs and monitoring as the core controls for agentic systems \cite{ref9}, \cite{ref10}, and empirical code-review research shows that reviewers need context to understand changes \cite{ref14}.

What is missing is the engineering synthesis of these findings into a single developer-facing pipeline. The agent platforms provide containers for capability, not for trust: none of them substitutes dummy values for the developer's real secrets, keeps the repository-host token out of the agent's reach, builds the diff in a component the agent could not tamper with, or holds the change behind an approval state machine before a PR is opened. The security studies characterise threats without delivering operational containment, and the governance papers specify controls without implementing them. VaultShield is positioned precisely in this gap. It combines Docker-based isolation, shape-preserving credential substitution, monitored execution with an append-only event log, optional egress allowlisting, trusted change capture and a human approval gate into one reproducible workflow, and it accompanies that workflow with an explicit statement of which risks remain. The following chapter describes how that synthesis was designed.

\clearpage

\chapter{Project Design}\label{project-design}

\section{Existing System}\label{existing-system}

In current practice, developers run AI coding agents in one of three ways. The most common is raw local execution: the agent's CLI or editor extension is started in the developer's own checkout, with the developer's shell environment, file-system permissions and network access. The agent reads and writes the working tree directly, executes shell commands as the developer, and its changes appear as uncommitted modifications that the developer may or may not inspect before committing. Some agent products add a permission layer on top, prompting before a command runs or honouring per-tool deny rules, but these layers are product-specific, configured in different files for different surfaces, and frequently relaxed by users who find repeated prompts tedious.

The second practice is ad-hoc containerisation. A developer who is aware of the risk may start the agent inside a container, typically by bind-mounting the working directory and sometimes the home directory or the Docker socket so that the agent's tooling continues to work. This improves isolation of the rest of the file system, but the mounted working tree is still the real one, so every write lands directly on the host. The real .env file travels into the container along with the source, and the container usually has unrestricted outbound network access. The third practice, development containers and cloud workspaces, provides a more reproducible environment but is designed for developer convenience rather than for constraining an untrusted actor; it does not by itself produce evidence of what an agent did or hold its output for review.

Figure 3.1 summarises the existing workflow. Its essential property is that execution and acceptance are a single step: by the time the developer looks at the result, the agent has already acted on the real repository with real credentials.


\begin{figure}[!htbp]
\centering
\SafeIncludeGraphics[width=\textwidth,height=0.62\textheight,keepaspectratio]{images/diagrams/fig3-1-existing-workflow.png}
\par\medskip
\captionof{figure}{Existing workflow – direct agent execution on the developer’s host}\label{fig:3-1}
\end{figure}

As shown in Figure 3.1, every arrow out of the agent touches a real asset. The limitations of this arrangement can be stated precisely. First, there is no isolation of side effects: file writes, deletions and command executions occur on the host and cannot be discarded wholesale. Second, there is no credential separation: the agent can read application secrets in .env, environment variables, cloud credential files and SSH keys. Third, there is no reliable evidence: terminal scroll-back is the only record of what happened, and it is neither structured nor durable. Fourth, there is no enforced review: the agent's changes sit in the working tree with nothing to stop them being committed and pushed unread. Fifth, cleanup is manual: a development server or background process started by the agent may keep running after the session ends. Finally, network activity is invisible: there is no record of which hosts the agent contacted.

\section{Proposed System}\label{proposed-system}

VaultShield replaces the single step of Figure 3.1 with a pipeline in which execution, evidence collection, storage and review are deliberately separate stages. The developer invokes vaultshield run with an agent, a prompt and a repository URL. The CLI loads the project configuration, generates a sortable, collision-resistant run identifier, creates a per-run directory and registers the run with the backend before any container exists, so that any later failure can be recorded against it. It then writes a dummy credential file that mirrors the shape of the developer's .env and resolves the agent adapter, refusing to proceed if the invocation is invalid.

Only then does the execution plane come into existence. The CLI creates a Docker container, bind-mounts only the dummy credential file, clones the repository inside the container and installs dependencies. If egress control is enabled, the sandbox is moved onto an internal network whose only peer is a filtering proxy. Before any agent code runs, the CLI records the pre-agent commit and Git objects through the Docker API as a trusted base. The agent then runs inside the container under a pseudo-terminal while an in-container watcher records file changes; both streams are merged into a timestamped event log. If routes are configured, the application is started inside the sandbox and screenshotted with a headless browser. The sandbox is then killed, and a separate hardened capture helper rebuilds the repository from the trusted base plus a filtered archive of the sandbox's work tree, commits the agent's changes to vaultshield/run-<run-id>, computes the structured diff and exports the branch as a Git bundle. Run outputs are scanned for the agent's own credential bytes, the sandbox is torn down in a finally block, and the artifacts are uploaded. The run moves to awaiting\_\allowbreak{}review.

A reviewer then opens the dashboard, inspects the diff, timeline, screenshots and automated security findings, and approves or rejects the run. Only an approved run can be turned into a pull or merge request, by running vaultshield submit <run-id> from a checkout of the same repository; a rejected run's branch is removed with vaultshield teardown <run-id>. Figure 3.2 presents this pipeline at a high level.


\begin{figure}[!htbp]
\centering
\SafeIncludeGraphics[width=\textwidth,height=0.62\textheight,keepaspectratio]{images/diagrams/fig3-2-proposed-pipeline.png}
\par\medskip
\captionof{figure}{Proposed VaultShield pipeline}\label{fig:3-2}
\end{figure}

As Figure 3.2 shows, the agent never touches the developer's checkout, and the path to a shared branch passes through a human decision. Table 3.1 contrasts the proposed system with the existing practices described earlier.

\Needspace{8\baselineskip}
\par\vspace{\TableCapGap}
\noindent\begin{minipage}{\textwidth}
\captionof{table}{Comparison of the existing system and the proposed system}\label{tab:3-1}
\end{minipage}
\par
\begin{longtable}[]{@{}|
  >{\raggedright\arraybackslash}p{(\columnwidth - 7\tabcolsep) * \real{0.3333}} |
  >{\raggedright\arraybackslash}p{(\columnwidth - 7\tabcolsep) * \real{0.3333}} |
  >{\raggedright\arraybackslash}p{(\columnwidth - 7\tabcolsep) * \real{0.3333}}|@{}}
\hline
Concern & Existing practice & Proposed system \\
\hline
\endfirsthead
\hline
Concern & Existing practice & Proposed system \\
\hline
\endhead
\endlastfoot
Working tree & Real checkout, modified in place & Fresh clone inside the container; host tree never mounted \\\hline
Application secrets & Real .env readable by the agent & Shape-preserving dummy values in a read-only .env.sandbox \\\hline
Repository token & Often in environment, remote URL or credential store visible to the agent & Passed only as the environment of the single git clone exec via a one-shot credential helper \\\hline
Network & Unrestricted, unlogged & Optional log-only or allowlist policy through a hardened proxy on an internal network, with every decision logged \\\hline
Evidence & Terminal scroll-back & JSON Lines event log, egress log, screenshots and structured diff in a per-run directory and backend \\\hline
Change capture & Uncommitted edits in the working tree & One commit on vaultshield/run-<id>, built by a separate helper from a trusted base, exported as a verified bundle \\\hline
Review & Optional and informal & Enforced state machine; PR/MR only from approved \\\hline
Cleanup & Manual & finally-guaranteed teardown plus a scheduled server-side reaper \\\hline
Security model & Usually implicit & Explicit list of implemented controls and residual risks \\\hline\end{longtable}
\addtocounter{table}{-1}
\par\medskip

The advantages follow directly from Table 3.1. Mistakes are contained, because the container and its clone are discarded after every run. Real secrets are not exposed, because the sandbox holds only dummy values and the repository token's lifetime is limited to one process. Every run produces durable, structured evidence identified by a single run identifier that also names the sandbox, the branch, the run directory and the backend record. The diff is trustworthy, because it is produced by tooling the agent could not influence. Review is not optional, because the only path to a PR passes through an approved status. Finally, the system is honest about what it does not protect against, which allows teams to layer stronger controls where required.

\section{Critical Components of System Architecture}\label{critical-components-of-system-architecture}

The architecture comprises three cooperating planes: a host-side CLI pipeline, a per-run Docker execution plane, and a backend and dashboard review plane. This section describes the five critical components required by the design, and a sixth group of supporting modules, in terms of their inputs, outputs and failure handling. Table 3.3 at the end of the section summarises them.

\subsection{Sandbox and Container Manager}\label{sandbox-and-container-manager}

The sandbox manager owns exactly one Docker container per run. Its inputs are the run identifier, the repository URL, the base image, the path of the generated .env.sandbox, and optional egress parameters. Sandbox.create ensures the image is present, verifies that the credential file exists, creates a container named vaultshield-run-<run-id> that is kept alive by an idle process, mounts the credential file read-only inside the container, clones the repository into the container's workspace and installs dependencies. The agent is later executed through exec primitives rather than as the container's main process, so that the container's lifetime bounds every process started inside it.

Private repositories are cloned without exposing the repository-host token to the agent. The token is resolved on the host and passed only as the environment of the single git clone exec, which Docker applies to that process alone; the command line carries only -c credential.helper=... naming the environment variables to read. The token therefore never appears in the container's environment, in the clone's Git configuration, in docker inspect output or in any later process. It is sent only to hosts that match a strict per-provider allowlist, so a look-alike host such as github.com.evil.com fails closed and is cloned unauthenticated. The Docker socket is never mounted into the sandbox.

When egress control is enabled, the companion EgressGuard builds the per-run network topology shown in Figure 3.3. Before the sandbox exists, it creates an ordinary bridge, an internal: true bridge with no default route and no gateway address, and a hardened proxy container attached to both. The sandbox is created on the open network so that cloning and dependency installation work, then switched onto the internal network only before any agent code runs. In these modes the sandbox also receives cap\_drop=["NET\_RAW"] and has IPv6 disabled, while the proxy runs with all capabilities dropped except NET\_\allowbreak{}BIND\_\allowbreak{}SERVICE, no-new-privileges, a read-only root file system, a 128 MiB memory limit, a 128-process limit and IP forwarding disabled.


\begin{figure}[!htbp]
\centering
\SafeIncludeGraphics[width=\textwidth,height=0.5\textheight,keepaspectratio]{images/diagrams/fig3-3-egress-topology.png}
\par\medskip
\captionof{figure}{Per-run egress network topology}\label{fig:3-3}
\end{figure}

As illustrated in Figure 3.3, the sandbox's only reachable peer after isolation is the proxy. A root process inside the sandbox cannot undo this, because changing routes requires CAP\_\allowbreak{}NET\_\allowbreak{}ADMIN, attaching networks requires the Docker API, raw packets require NET\_RAW and the proxy refuses to start if IP forwarding is enabled. The proxy is a custom standard-library CONNECT and TLS-SNI tunnelling proxy; it does not terminate TLS. It filters by destination host against the union of the agent's built-in allowlist and the project's egress.allow entries and denies IP literals, port mismatches and names that resolve to private addresses.

Failure handling is uniform. Any exception during creation is caught, the partially created container is removed, and the error is re-raised as a single SandboxError type with a message that never contains a token. teardown calls remove(v=True, force=True), which kills every process in the container and removes anonymous volumes; it is idempotent and tolerates an already-removed container. The run pipeline calls it from a finally block, so teardown executes on success, agent failure, readiness timeout, capture failure and keyboard interrupt alike, and the egress guard then removes the proxy and both networks on a best-effort basis.

\subsection{Run Monitoring and Logging}\label{run-monitoring-and-logging}

The monitoring component takes as input the started sandbox, the agent adapter's SandboxInvocation and the project's off\_\allowbreak{}limits\_\allowbreak{}paths. It produces the run's Event\_Log, events.jsonl, a JSON Lines file in which each line is an object with event\_type, timestamp and payload. Three event types exist: output for terminal output captured from the agent's pseudo-terminal, file\_\allowbreak{}change for changes observed by the in-container watcher, and violation for off-limits writes, egress denials and token-hygiene findings. JSON Lines was chosen because it can be appended incrementally and a partially written log remains parseable line by line, so even a crashed run leaves readable evidence.

The file-change watcher runs inside the sandbox so that it observes the agent's real writes rather than a host view. It is a small standard-library program that snapshots modification times and sizes across the cloned repository, polls for created, modified and deleted files, and on shutdown performs a final reconciliation pass against the original baseline to report net changes the polling loop may have missed. The host merges watcher records with PTY output into a single list sorted stably by timestamp, then appends it to the Event\_Log. Each changed path is matched against the normalised off\_\allowbreak{}limits\_\allowbreak{}paths entries by both directory containment and glob matching; the first match produces a violation event and the run fails with OffLimitsViolation, but only after the Event\_Log has been written, so the evidence always outlives the failure. The matcher deliberately errs towards false positives, because a false positive fails one visible run while a false negative silently permits a forbidden write.

Two further evidence producers complete the component. Browser capture starts the application inside the sandbox with the configured start\_\allowbreak{}command, polls ready\_path until it returns HTTP 200 within --ready-timeout seconds, and captures one PNG screenshot per configured route with Playwright and Chromium on the host. The two failure classes are kept distinct: ReadinessTimeout means the application never came up and yields zero screenshots, while ScreenshotError means the application answered but visual capture failed, and partial screenshots are kept. Token hygiene snapshots the agent's own declared credential files through the Docker API before and after the run, collects their secret bytes as ``needles'', and scans diff.json, events.jsonl and changed files for them. A hit fails the run with a TokenHygieneViolation whose message never contains the matched bytes, renames the diff to diff.quarantined.json, and causes uploaded event payloads to be scrubbed with ***.

\subsection{Diff Capture on a Dedicated Branch}\label{diff-capture-on-a-dedicated-branch}

The capture component addresses an integrity problem: if the branch and diff were built inside the agent's own sandbox, a hostile agent could tamper with Git state, hooks, filters or configuration to misrepresent its changes. VaultShield therefore splits capture across two moments and two containers. Before any agent code runs, record\_\allowbreak{}trusted\_\allowbreak{}base reads the pre-agent commit and a filtered snapshot of the Git object store through the Docker API. After the agent finishes, kill\_\allowbreak{}sandbox stops every sandbox process and archive\_\allowbreak{}sandbox\_\allowbreak{}repo extracts a filtered tar archive of the work tree; the filter fails closed on malformed headers, path traversal, symlinked parents, nested .git directories and oversized streams, and normalises member metadata.

A fresh capture helper container is then created from the same image ID as the sandbox but with no network, no capabilities, no volumes and no environment. It lays down the filtered work tree, initialises a new repository from the trusted base objects, refuses the capture if the untrusted Git configuration defines filter, textconv or diff drivers, and checks that .git is a plain directory at the repository root rather than a gitfile or symlink. commit\_\allowbreak{}run\_\allowbreak{}branch creates vaultshield/run-<run-id>, stages every change with git add --all, and commits with --allow-empty --no-verify under a fixed committer identity, so that exactly one commit is created even if the agent changed nothing. build\_\allowbreak{}structured\_\allowbreak{}diff runs git diff --numstat -z --no-renames --no-ext-diff --no-textconv between the base and the branch, then a per-file git diff -U0 for text files, and writes diff.json as a list of \{path, additions, deletions, patch, patch\_\allowbreak{}truncated\} entries.

The diff obeys explicit size limits so that it can always be uploaded: each patch is capped at 200 KiB and cut at a line boundary, and the whole request body is capped at 2 MiB, below the backend's 2.5 MiB upload limit. Once the budget is spent, later text files carry "patch": "" with "patch\_\allowbreak{}truncated": true, so the reviewer and the security scan can see exactly where coverage ends. If per-file metadata alone exceeds the cap, capture fails with a clear error rather than silently dropping entries. After token hygiene passes, the branch is exported as branch.bundle together with branch.json, because the helper is destroyed at the end of the run and the branch never existed in the developer's checkout. The outputs of this component are therefore the branch, diff.json, and the bundle metadata; its failures fail the run, record a violation event where appropriate, and never produce a bundle for a quarantined diff.

\subsection{Review and Human Approval Service}\label{review-and-human-approval-service}

The review service spans the backend state machine, the dashboard decision controls and the CLI's submit and teardown commands. Its input is a fully uploaded run in awaiting\_\allowbreak{}review; its outputs are a recorded Approval and, for approved runs, a PR or MR URL. The Run.status field is a six-value state machine, shown in Figure 3.4. The CLI may PATCH a run from running to awaiting\_\allowbreak{}review or to failed; both transitions are gated to originate only from running.


\begin{figure}[!htbp]
\centering
\SafeIncludeGraphics[width=\textwidth,height=0.7\textheight,keepaspectratio]{images/diagrams/fig3-4-run-state-machine.png}
\par\medskip
\captionof{figure}{Run status state machine}\label{fig:3-4}
\end{figure}

The approve and reject endpoints call apply\_\allowbreak{}decision, which performs the gate as a single conditional update inside a MongoDB transaction: the run's status changes only if it is currently awaiting\_\allowbreak{}review, and an Approval row is created in the same transaction. Because Approval.run is a one-to-one field, a run can be decided at most once at the schema level. Of two concurrent decisions, MongoDB's single-document atomicity guarantees that exactly one matches the condition; the loser's transaction is retried on TransientTransactionError and then finds the winner's decision, returning HTTP 409 with the run's actual status rather than a server error. The dashboard shows the Approve and Reject buttons only for runs in awaiting\_\allowbreak{}review, disables both while a decision is in flight, and after approval displays the copy ``approved, awaiting submission'' together with the exact vaultshield submit <run-id> command.

vaultshield submit reads the run from the backend and refuses, with the run's actual status in the message, unless the status is approved; this check precedes every side effect, so a refused submission performs one GET and nothing else. For an approved run it verifies that the current checkout's origin is the repository the run was taken from, imports the branch from branch.bundle after verifying its SHA-256, resolves the developer's Git host credential, performs a plain git push without --force, opens the PR or MR through the GitHub or GitLab API with a description identifying the run and its review, and records the pr\_url on the run. vaultshield teardown on a rejected run deletes the run branch from the remote if it was pushed and deletes the local imported branch only if it still points at the recorded tip.

\subsection{Backend API and Dashboard}\label{backend-api-and-dashboard}

The backend is a Django 5.2 project with Django REST Framework, using MongoDB through django-mongodb-backend as its system of record. Its inputs are authenticated HTTP requests from the CLI and the dashboard; its outputs are JSON representations of projects and runs and stored artifacts. Token authentication and the IsAuthenticated permission are configured globally, so every /api/ path except login requires an Authorization: Token <token> header. The REST surface is summarised in Table 3.2. Runs are addressed by the CLI-generated run\_id rather than the database primary key, and there are deliberately no PUT routes and no run DELETE route.

\Needspace{8\baselineskip}
\par\vspace{\TableCapGap}
\noindent\begin{minipage}{\textwidth}
\captionof{table}{Backend REST API surface}\label{tab:3-2}
\end{minipage}
\par
\begin{longtable}[]{@{}|
  >{\raggedright\arraybackslash}p{(\columnwidth - 7\tabcolsep) * \real{0.3333}} |
  >{\raggedright\arraybackslash}p{(\columnwidth - 7\tabcolsep) * \real{0.3333}} |
  >{\raggedright\arraybackslash}p{(\columnwidth - 7\tabcolsep) * \real{0.3333}}|@{}}
\hline
Method & Path & Purpose \\
\hline
\endfirsthead
\hline
Method & Path & Purpose \\
\hline
\endhead
\endlastfoot
POST & /api/auth/login/ & Exchange username and password for a dashboard token \\\hline
GET, POST & /api/projects/ & List or create projects \\\hline
DELETE & /api/projects/:id/ & Delete a project and its runs and artifacts \\\hline
POST & /api/runs/ & Register a run \\\hline
GET & /api/runs/?status=\allowbreak{}\&submitted=\allowbreak{}\&project= & Filtered, paginated run list \\\hline
GET & /api/runs/:run\_id/ & Full run detail: events, diff, screenshots and security findings \\\hline
PATCH & /api/runs/:run\_id/ & Gated status transitions and pr\_url \\\hline
POST & /api/runs/:run\_\allowbreak{}id/events/ & Bulk Event\_Log ingest with order preserved by sequence \\\hline
POST & /api/runs/:run\_\allowbreak{}id/diff/ & Store the structured diff and enqueue the security scan \\\hline
POST & /api/runs/:run\_\allowbreak{}id/screenshots/ & Upload one screenshot \\\hline
POST & /api/runs/:run\_\allowbreak{}id/approve/, /reject/ & Record the review decision \\\hline\end{longtable}
\addtocounter{table}{-1}
\par\medskip

Several backend mechanisms are specific to MongoDB. Event sequence numbers are allocated with a single atomic find\_\allowbreak{}one\_\allowbreak{}and\_\allowbreak{}update using \$inc on the run document, replacing the row lock that a relational database would use, and a unique compound index on turns any regression into a loud integrity error. Multi-document transactions, required for the approval gate, event batches and the security scan's delete-then-insert, are available only on a replica set, so the backend requires MongoDB Atlas or a replica-set deployment. Asynchronous work runs on Celery with Redis as broker: diff ingest enqueues a static security scan over added lines only, which never executes diff content; Celery Beat runs a sandbox reaper every 15 minutes and flags approved-but-unsubmitted runs as stuck every hour. The local deployment is a Docker Compose stack of six services: redis, web, channels, celery-worker, celery-beat and frontend.

The dashboard is a React 19 single-page application styled with Tailwind CSS 4 and built with Vite 7. Its API client is a hand-written fetch wrapper that attaches the token to every call except login and redirects to the sign-in screen on HTTP 401. The Vite server proxies /api and /media to the backend so that the browser remains same-origin and the backend needs no CORS middleware. The interface consists of a landing page, a sign-in page, a projects overview, a run history table with status filters and search, and a run detail page with four tabs and the decision controls. Failures in this component are surfaced explicitly: a refused decision shows the status the backend reports, a network failure states that no decision was recorded, and a protection-log run uploaded by the protect sync command shows ``Protection log: nothing to review'' instead of decision buttons.

\subsection{Supporting Modules: Protect and Orchestrate}\label{supporting-modules-protect-and-orchestrate}

Two supporting command groups complement the core pipeline. Protect addresses agents that developers continue to run outside the sandbox: it moves every non-empty value in .env and .env.* files into the macOS Keychain, replaces each value in the file with the placeholder vs\_\allowbreak{}protected, and provides vaultshield exec -- <command> to run the application with the real values injected into its environment. It can also install agent-specific deny rules and a PreToolUse hook for Antigravity, Claude Code, Kiro and Codex as backstops. Its design follows directly from the spike results, where keeping the real value off disk was the only layer that held on every surface. Orchestrate scans the current project, asks a few clarifying questions through a Gemini model and produces an enriched agent prompt; a final secret check guarantees that the prompt contains no secret value or secret-file content before the run is registered.

\Needspace{8\baselineskip}
\par\vspace{\TableCapGap}
\noindent\begin{minipage}{\textwidth}
\captionof{table}{Summary of critical components -- inputs, outputs and failure handling}\label{tab:3-3}
\end{minipage}
\par
\begin{longtable}[]{@{}|
  >{\raggedright\arraybackslash}p{(\columnwidth - 9\tabcolsep) * \real{0.2500}} |
  >{\raggedright\arraybackslash}p{(\columnwidth - 9\tabcolsep) * \real{0.2500}} |
  >{\raggedright\arraybackslash}p{(\columnwidth - 9\tabcolsep) * \real{0.2500}} |
  >{\raggedright\arraybackslash}p{(\columnwidth - 9\tabcolsep) * \real{0.2500}}|@{}}
\hline
Component & Main inputs & Main outputs & Failure handling \\
\hline
\endfirsthead
\hline
Component & Main inputs & Main outputs & Failure handling \\
\hline
\endhead
\endlastfoot
Sandbox and container manager & Run ID, repo URL, base image, .env.sandbox, egress policy & Running container with the cloned repository; optional internal network and proxy & Partial container removed on any creation error; single SandboxError; teardown in finally; best-effort network and proxy removal \\\hline
Run monitoring and logging & Sandbox, SandboxInvocation, off\_\allowbreak{}limits\_\allowbreak{}paths, routes & events.jsonl, screenshots, violation events & Log written before failure is raised; distinct OffLimitsViolation, ReadinessTimeout, ScreenshotError, TokenHygieneViolation \\\hline
Diff capture on a dedicated branch & Trusted base, filtered work-tree archive & vaultshield/run-<id> commit, diff.json, branch.bundle, branch.json & Fail-closed archive filter; UnsafeGitConfig refusal; size-capped patches; quarantine on hygiene hit \\\hline
Review and human approval service & Run in awaiting\_\allowbreak{}review, reviewer decision & Approval record; PR/MR URL; branch cleanup & Conditional update in a transaction; 409 on conflict; submit refuses non-approved runs before any side effect \\\hline
Backend API and dashboard & Authenticated REST calls & Persisted runs and artifacts; rendered evidence & Global token authentication; validated ingest; Celery reaper and stuck-approval flag \\\hline
Supporting modules & Host .env files; task description & Keychain-held values and placeholders; enriched prompt & Whole-file refusal on unparseable lines with rollback; final secret check on prompts \\\hline\end{longtable}
\addtocounter{table}{-1}
\par\medskip

\begin{figure}[H]
\centering
\SafeIncludeGraphics[width=\textwidth,height=0.86\textheight,keepaspectratio]{images/diagrams/fig3-5-system-architecture.png}
\par\medskip
\captionof{figure}{Overall system architecture}\label{fig:3-5}
\end{figure}

Figure 3.5 shows all components and the data flows between them. The CLI on the developer's host is the only component that talks to the Docker daemon; it creates the sandbox, streams the agent's output, reads the trusted base and the work-tree archive through the Docker API, and drives the capture helper. The dummy credential file is the only host path mounted into the sandbox, and the sandbox reaches model APIs either directly or only through the proxy. Artifacts flow from the run directory to the backend over authenticated HTTP; the backend stores them in MongoDB and the media volume and hands security scanning to Celery through Redis. The reviewer interacts only with the dashboard, and the developer's submit command is the only component that pushes to the repository host, after confirming approval with the backend.

\section{System Diagrams}\label{system-diagrams}

System diagrams give a visual representation of the structure, behaviour and interactions of the proposed VaultShield system. They show the main components of the system and how information flows between the developer, the reviewer, the sandbox and the backend. The diagrams used for the proposed system are the UML Diagram, Activity Diagram, Use Case Diagram and Sequence Diagram.

\subsection{UML Diagram}\label{uml-diagram}

Figure 3.6 shows the persistent backend entities, Project, Run and the records attached to a run, together with the principal CLI-side classes AgentAdapter, EgressGuard, Sandbox and TrustedCapture. Run is the hub of the schema. RunEvent, Screenshot and SecurityFinding are many-per-run, as is the SandboxRecord used by the reaper, which is omitted from the diagram for clarity, while RunDiff and Approval are one-to-one, which is what makes ``a run has exactly one change set and at most one decision'' a schema-level guarantee; every relation cascades on delete. In the terminology of the problem statement, a review request is a Run in awaiting\_\allowbreak{}review, and its resolution is the Approval record. Notably, no backend field can store a repository-host token; pr\_url, a public URL, is the only value submit writes back. The MVP's User, also omitted, authenticates with a token but is not linked to projects or approvals, reflecting the single-developer scope. On the CLI side, an AgentAdapter describes how its agent runs inside the Sandbox rather than spawning a host process, the EgressGuard isolates the Sandbox's network, monitoring of the Sandbox produces the RunEvent records, and the TrustedCapture helper produces the file entries stored in RunDiff.

\begin{figure}[H]
\centering
\SafeIncludeGraphics[width=\textwidth,height=0.85\textheight,keepaspectratio]{images/diagrams/fig3-6-class.png}
\par\medskip
\captionof{figure}{UML class diagram of the core backend entities and CLI classes}\label{fig:3-6}
\end{figure}

\subsection{Activity Diagram}\label{activity-diagram}

Figure 3.7 traces a run from setup through review and submission, reading each row in the direction of its arrows. The ordering enforces key guarantees: errors are registered before sandbox creation, invalid agent configurations never start a container, policy failures are recorded, and teardown runs on every exit path. Reviewer approval is the only route to submission, and the CLI checks approval again immediately before pushing.

\begin{figure}[H]
\centering
\SafeIncludeGraphics[width=\textwidth,height=0.88\textheight,keepaspectratio]{images/diagrams/fig3-7-activity.png}
\par\medskip
\captionof{figure}{Activity diagram of a VaultShield run, from agent kickoff to PR creation}\label{fig:3-7}
\end{figure}

\subsection{Use Case Diagram}\label{use-case-diagram}

Figure 3.8 shows the principal VaultShield actors and their interactions. The Developer configures and executes sandboxed runs, then submits or removes them. The Reviewer/Approver inspects run evidence and records an approval or rejection; in the single-developer MVP, one person may perform both roles, but they remain distinct workflow roles. The System actor represents automated scanning and cleanup by Celery workers. Executing a run includes dummy-credential generation, monitoring, trusted capture and upload, and may be extended by egress enforcement and screenshot capture. Submission is permitted only after approval, and the Repository host actor receives the pushed branch and pull request.

\begin{figure}[H]
\centering
\SafeIncludeGraphics[width=\textwidth,height=0.72\textheight,keepaspectratio]{images/diagrams/fig3-8-use-case-diagram.png}
\par\medskip
\captionof{figure}{Use case diagram}\label{fig:3-8}
\end{figure}

\subsection{Sequence Diagram}\label{sequence-diagram}

\begin{figure}[H]
\centering
\SafeIncludeGraphics[width=\textwidth,height=0.8\textheight,keepaspectratio]{images/diagrams/fig3-9-sequence.png}
\par\medskip
\captionof{figure}{Sequence diagram of one complete approved run}\label{fig:3-9}
\end{figure}

Figure 3.9 shows the message exchange for a single successful, approved run. Messages 1 to 9 cover the run itself, 10 and 11 the review, and 12 to 16 the submission. Three properties of the sequence are significant. First, the backend learns of the run before Docker is asked to create anything, which is what allows every subsequent failure to be recorded. Second, the agent's activity is observed from inside the container, but the branch and diff are produced by a different container from data read through the Docker API, so the evidence the reviewer sees is not generated by tooling the agent controlled. Third, the approval is a single transactional write, and submit re-reads the status from the backend immediately before pushing, so the PR path cannot be taken from a stale client view.

\clearpage

\chapter{Project Planning and Cost Estimation}\label{project-planning}

\section{Project Development Methodology}\label{project-development-methodology}

VaultShield was developed using an iterative and incremental methodology with specification-driven development, a combination of agile practice and lightweight formal specification. Before implementation, the team wrote a requirements specification in which every requirement carried numbered acceptance criteria, and a design document that stated the architecture, data model and a set of numbered correctness properties. The design was then decomposed into small, independently testable tasks. These identifiers are cited throughout the source code and test suite, which provides traceability from each requirement to the code and test that satisfy it.

Implementation proceeded in short iterations. Each iteration selected a coherent slice of tasks, implemented them, added unit and property-based tests asserting the relevant acceptance criteria and correctness properties, and ended with a documentation update and, where the slice touched isolation, a Docker-backed integration test. Larger feature groups were delivered as numbered batches. The approach suited the project for two reasons: the security-sensitive behaviour could be specified precisely in advance and verified as properties, and the scope could be extended iteratively without destabilising earlier work, because the existing tests guarded every prior guarantee. The major phases of the development methodology are described below:

\begin{enumerate}
\item
  \textbf{Problem Study and Literature Survey:} Study coding agents, their risks, container isolation and code review, and identify the research gap that the project addresses.
\item
  \textbf{Feasibility Spike:} Measure how a production agent handles .env secrets under four protection configurations, and probe container isolation.
\item
  \textbf{Requirements Specification:} Write numbered functional and security requirements, each with acceptance criteria.
\item
  \textbf{System Design:} Define the architecture, run lifecycle, data model, REST API, threat model and correctness properties.
\item
  \textbf{CLI Core:} Implement configuration, the run directory, credential shaping, the sandbox lifecycle and the agent adapters.
\item
  \textbf{Evidence Capture:} Implement PTY and file-change monitoring, off-limits enforcement, screenshots and Git capture.
\item
  \textbf{Backend:} Implement the models, REST API, authentication, artifact ingest, security scan and scheduled tasks.
\item
  \textbf{Dashboard:} Implement the project list, run history, run detail tabs and approval controls.
\item
  \textbf{Submission:} Implement PR/MR submission, teardown, the GitHub and GitLab providers, and branch bundles.
\item
  \textbf{Hardening:} Add the egress topology and proxy, the trusted capture helper and token hygiene.
\item
  \textbf{Datastore Migration:} Move the backend from PostgreSQL to MongoDB, with transactional approval and atomic sequence allocation.
\item
  \textbf{Supporting Modules:} Add the Protect and Orchestrate modules and the one-line installer.
\item
  \textbf{User Interface:} Add routing, the redesign, accessibility checks, responsive layout and the theme switch.
\item
  \textbf{Evaluation and Documentation:} Run the tests and performance benchmark and write this report.
\end{enumerate}

\section{Project Schedule and Gantt Chart}\label{project-schedule}

The project schedule defines the major activities involved in the development of VaultShield and their planned sequence. The Semester VII activities follow the project Gantt chart, which is reproduced in full as Figure~5.21; the Semester VIII activities are planned extensions of the work. Table 4.1 lists each activity with its duration and target period.

\begin{table}[H]
\centering\small
\caption{Planned project schedule for VaultShield}\label{tab:4-2}
\begin{tabularx}{\textwidth}{|>{\centering\arraybackslash}p{0.35in}|>{\raggedright\arraybackslash}X|>{\raggedright\arraybackslash}p{0.85in}|>{\raggedright\arraybackslash}p{1.55in}|}\hline
\textbf{Sr.} & \textbf{Project activity} & \textbf{Duration} & \textbf{Target period} \\\hline
\multicolumn{4}{|l|}{\textbf{Semester VII}} \\\hline
1 & Group formation and topic finalisation; identifying the scope and objectives of the project & 1 week & 10 Jul -- 17 Jul 2026 \\\hline
2 & Discussing the project topic with the help of a paper prototype & 1 week & 17 Jul -- 24 Jul 2026 \\\hline
3 & Identifying the functionalities of the project as per the objectives identified & 1 week & 24 Jul -- 31 Jul 2026 \\\hline
4 & Literature review and research paper study & 1 week & 31 Jul -- 7 Aug 2026 \\\hline
5 & Designing the graphical user interface & 3 weeks & 7 Aug -- 28 Aug 2026 \\\hline
6 & Presentation I & 1 day & 4 Sep 2026 \\\hline
7 & Policy logic development & 1 week & 18 Sep -- 25 Sep 2026 \\\hline
8 & Core model development and training & 1 week & 25 Sep -- 1 Oct 2026 \\\hline
9 & Presentation II & 1 day & 9 Oct 2026 \\\hline
10 & Incorporating the suggestions from Presentation II and submitting a draft copy to the project guide for proofreading & 1 week & 16 Oct -- 23 Oct 2026 \\\hline
11 & Draft submission to the guide for review & 2 days & 23 Oct -- 24 Oct 2026 \\\hline
12 & Final black book submission & 1 week & 26 Oct -- 30 Oct 2026 \\\hline
13 & External presentation and viva voce & 1 day & 2 Nov 2026 \\\hline
\multicolumn{4}{|l|}{\textbf{Semester VIII}} \\\hline
14 & Evaluation with real agents on sample repositories & 3 weeks & January 2027 \\\hline
15 & Default sandbox hardening and fixes for issues found in testing & 3 weeks & January -- February 2027 \\\hline
16 & Verification of the Codex and Kiro egress allowlists; reviewer-identity separation study & 2 weeks & February 2027 \\\hline
17 & Usability study with reviewers and dashboard refinements & 2 weeks & March 2027 \\\hline
18 & Research paper drafting and submission & 2 weeks & March 2027 \\\hline
19 & Final report, demonstration and viva preparation & 3 weeks & March -- April 2027 \\\hline
\end{tabularx}
\end{table}

The schedule provides an overview of the planned development activities. The Semester VIII dates are targets and may change depending on the evaluation results, the issues found during hardening, and the academic calendar.

\section{Team Roles and Responsibilities}\label{team-roles-and-responsibilities}

The team of four divided responsibilities by architectural plane, with shared ownership of testing and documentation. Table 4.2 records the allocation.

\Needspace{8\baselineskip}
\par\vspace{\TableCapGap}
\noindent\begin{minipage}{\textwidth}
\captionof{table}{Team roles and responsibilities}\label{tab:4-3}
\end{minipage}
\par
\begin{longtable}[]{@{}|
  >{\raggedright\arraybackslash}p{(\columnwidth - 7\tabcolsep) * \real{0.3333}} |
  >{\raggedright\arraybackslash}p{(\columnwidth - 7\tabcolsep) * \real{0.3333}} |
  >{\raggedright\arraybackslash}p{(\columnwidth - 7\tabcolsep) * \real{0.3333}}|@{}}
\hline
Member & Primary role & Responsibilities \\
\hline
\endfirsthead
\hline
Member & Primary role & Responsibilities \\
\hline
\endhead
\endlastfoot
Aayush Subhash Balip & Backend and data engineering & Django/DRF API, MongoDB models and migration, run state machine and approval atomicity, artifact ingest, Celery security scan and scheduled tasks, Docker Compose stack \\\hline
Krish Mahesh Patil & Frontend and user experience & React dashboard, API client, run detail tabs and decision controls, routing, accessibility, responsive layout, landing page \\\hline
Muhammad Anis Momin & Security research and sandbox engineering & Sandbox lifecycle, credential shaping, egress proxy and topology, trusted capture, token hygiene, threat model, isolation probes and the secret-exposure spike \\\hline
Sameer Hasan Saiyed & Integration, testing and documentation & CLI pipeline integration, agent adapters, submit and teardown, test strategy and property-based tests, benchmark, README and report \\\hline\end{longtable}
\addtocounter{table}{-1}
\par\medskip

\section{Software Size Estimation}\label{preliminary-software-size-estimation-kloc}

Software size estimation approximates the size of the software to be developed and provides the input for estimating development effort and schedule. For VaultShield, the size is expressed in thousand lines of code, measured from the delivered source rather than assumed.

\subsection{KLOC-Based Estimation}\label{kloc-based-estimation}

The size of the delivered software was measured directly from the repository rather than estimated, using logical source lines: non-blank lines that are not comments and, for Python, not docstrings. This measure was chosen because the code base is heavily documented, and counting physical lines would overstate the size of the executable logic. Migrations, lock files and generated content were excluded. Table 4.3 reports the counts.

\Needspace{8\baselineskip}
\par\vspace{\TableCapGap}
\noindent\begin{minipage}{\textwidth}
\captionof{table}{Measured software size by component}\label{tab:4-4}
\end{minipage}
\par
\begin{longtable}[]{@{}|
  >{\raggedright\arraybackslash}p{(\columnwidth - 7\tabcolsep) * \real{0.2727}} |
  >{\raggedleft\arraybackslash}p{(\columnwidth - 7\tabcolsep) * \real{0.3636}} |
  >{\raggedleft\arraybackslash}p{(\columnwidth - 7\tabcolsep) * \real{0.3636}}|@{}}
\hline
Component & Product code, SLOC & Test code, SLOC \\
\hline
\endfirsthead
\hline
Component & Product code, SLOC & Test code, SLOC \\
\hline
\endhead
\endlastfoot
CLI package & 11,954 & 16,697 \\\hline
CLI tools & 450 & -- \\\hline
Backend & 1,150 & 2,552 \\\hline
Frontend & 4,356 & 2,511 \\\hline
Deployment configuration & 160 & -- \\\hline
\textbf{Total delivered product} & \textbf{18,070} & \textbf{21,760} \\\hline
Research spike scripts & 1,642 & -- \\\hline\end{longtable}
\addtocounter{table}{-1}
\par\medskip

The delivered product therefore measures KLOC = 18.07. The test code is larger than the product code, which reflects the emphasis on verifying security properties; it is excluded from the primary estimate, consistent with the convention that COCOMO sizes delivered source instructions, and is considered separately as a sensitivity case.

\section{Effort Estimation Using COCOMO}\label{effort-estimation-using-basic-cocomo}

\subsection{Basic COCOMO Model}\label{basic-cocomo-model}

The Basic Constructive Cost Model \cite{ref15} estimates effort and schedule from size alone. VaultShield was classified as an organic project: it was built by a small team in a familiar environment, with flexible requirements and no hard real-time or hardware constraints. The organic-mode coefficients are a = 2.4, b = 1.05, c = 2.5 and d = 0.38, giving the following equations, where E is effort in person-months, T is development time in months and P is the average number of persons.

E = a × (KLOC)\^{}b = 2.4 × (KLOC)\^{}1.05

T = c × (E)\^{}d = 2.5 × (E)\^{}0.38

P = E / T

\subsection{Calculation}\label{cocomo-calculation}

Effort. With KLOC = 18.07, the exponent term is (18.07)\^{}1.05 = e\^{}(1.05 × ln 18.07) = e\^{}(1.05 × 2.8943) = e\^{}3.0390 = 20.884. Hence

E = 2.4 × 20.884 = 50.12 person-months.

Development time. The exponent term is (50.12)\^{}0.38 = e\^{}(0.38 × ln 50.12) = e\^{}(0.38 × 3.9144) = e\^{}1.4875 = 4.426. Hence

T = 2.5 × 4.426 = 11.06 months.

Average staffing. P = 50.12 / 11.06 = 4.53 persons.

\Needspace{8\baselineskip}
\par\vspace{\TableCapGap}
\noindent\begin{minipage}{\textwidth}
\captionof{table}{Basic COCOMO estimate}\label{tab:4-5}
\end{minipage}
\par
\begin{longtable}[]{@{}|
  >{\raggedright\arraybackslash}p{(\columnwidth - 9\tabcolsep) * \real{0.2143}} |
  >{\raggedright\arraybackslash}p{(\columnwidth - 9\tabcolsep) * \real{0.2143}} |
  >{\raggedleft\arraybackslash}p{(\columnwidth - 9\tabcolsep) * \real{0.2857}} |
  >{\raggedleft\arraybackslash}p{(\columnwidth - 9\tabcolsep) * \real{0.2857}}|@{}}
\hline
Quantity & Formula & Product code only, KLOC = 18.07 & Sensitivity: product + tests, KLOC = 39.83 \\
\hline
\endfirsthead
\hline
Quantity & Formula & Product code only, KLOC = 18.07 & Sensitivity: product + tests, KLOC = 39.83 \\
\hline
\endhead
\endlastfoot
Effort E & 2.4 × KLOC\^{}1.05 & 50.12 person-months & 114.93 person-months \\\hline
Development time T & 2.5 × E\^{}0.38 & 11.06 months & 15.17 months \\\hline
Average staffing P & E / T & 4.53 persons & 7.58 persons \\\hline\end{longtable}
\addtocounter{table}{-1}
\par\medskip

The primary estimate of about 4.5 persons over about 11 months agrees closely with the actual configuration of a four-member team working across two semesters, which supports the organic classification. The sensitivity case shows that if the test suite is treated as delivered code of equal difficulty, the model would call for roughly 7.6 persons over 15 months. Basic COCOMO was calibrated on projects from the 1970s and does not account for modern framework reuse or for current development tooling, so these figures are best read as a conventional-effort reference rather than as a measurement of the team's actual hours. 

\section{Human Resource Cost}\label{cost-estimation-human-resources}

The project was completed by the four student team members as part of their academic programme. No salary, stipend, contractor fee, consultant fee, or external reviewer payment was made for this project. Accordingly, the table below records actual cash expenditure rather than assigning a hypothetical monetary value to unpaid student effort. The COCOMO result is an effort and schedule estimate only; it is not treated as money spent.

\begin{table}[H]
\centering\small
\caption{Actual human-resource cash expenditure}\label{tab:4-6}
\begin{tabularx}{\textwidth}{|>{\raggedright\arraybackslash}X|>{\raggedright\arraybackslash}X|>{\raggedleft\arraybackslash}p{1.45in}|}\hline
\textbf{Cost item} & \textbf{Actual basis} & \textbf{Actual cost in \rupee} \\\hline
Student development effort & Unpaid work completed as part of the academic project & 0 \\\hline
Guide and co-guide supervision & Academic guidance; no project-specific payment & 0 \\\hline
Review and user acceptance & No paid external reviewer or consultant engaged & 0 \\\hline
\textbf{Total actual human-resource cash expenditure} & & \textbf{0} \\\hline
\end{tabularx}
\end{table}

\section{Hardware, Software and Cloud Infrastructure Cost}\label{cost-estimation-hardware-software-and-infrastructure}

The work used equipment already available to the team, open-source software, local development and testing infrastructure, and services available without an additional project-specific payment. No new hardware was purchased, no paid cloud environment was commissioned for the MVP, and no incremental usage fee was paid from the project budget. The table records the project-specific cash amount, not the market value or depreciation of resources already owned or available.

\begin{table}[H]
\centering\small
\caption{Actual hardware, software and infrastructure cash expenditure}\label{tab:4-7}
\begin{tabularx}{\textwidth}{|>{\raggedright\arraybackslash}p{1.18in}|>{\raggedright\arraybackslash}X|>{\raggedright\arraybackslash}p{1.55in}|>{\raggedleft\arraybackslash}p{0.95in}|}\hline
\textbf{Category} & \textbf{Item or resource} & \textbf{Actual basis} & \textbf{Cost in \rupee} \\\hline
Hardware & Existing team laptops and devices & No new equipment purchased & 0 \\\hline
Software & Python, Django/DRF, React, Git, Docker tooling, test frameworks and other project dependencies & Open-source software or tools already available at no additional project charge & 0 \\\hline
Repository and collaboration & Git repository hosting and collaboration tools & Existing or free-plan access & 0 \\\hline
Database and testing & Local MongoDB replica-set container and local test stack & Local/open-source testing; no paid test environment & 0 \\\hline
Connectivity and utilities & Existing internet and electricity arrangements & No incremental project-specific payment separately incurred & 0 \\\hline
Cloud and model/API services & No paid production cloud deployment or project-funded usage charge recorded & No additional project expenditure recorded & 0 \\\hline
\textbf{Total hardware, software and infrastructure cash expenditure} & & & \textbf{0} \\\hline
\end{tabularx}
\end{table}

\section{Testing, Deployment and Maintenance Cost}\label{cost-estimation-testing-deployment-and-maintenance}

Testing was performed using local machines and the local Docker-based environment. The MVP was not deployed to a paid cloud environment. Future maintenance may require effort, but no future or hypothetical maintenance amount is included in the actual-cost total because it has not been incurred as project expenditure.

\begin{table}[H]
\centering\small
\caption{Actual testing, deployment and maintenance cash expenditure}\label{tab:4-8}
\begin{tabularx}{\textwidth}{|>{\raggedright\arraybackslash}X|>{\raggedright\arraybackslash}X|>{\raggedleft\arraybackslash}p{1.25in}|}\hline
\textbf{Item} & \textbf{Actual basis} & \textbf{Actual cost in \rupee} \\\hline
Automated test execution & Local machines and open-source test frameworks & 0 \\\hline
Docker-backed integration tests and benchmarks & Local Docker-based test environment & 0 \\\hline
MVP deployment & Local Docker Compose stack; no paid cloud deployment & 0 \\\hline
Cloud VM or managed paid database & Not commissioned for this MVP & 0 \\\hline
Post-submission maintenance & Future work; no maintenance payment incurred during the project & 0 \\\hline
\textbf{Total actual testing, deployment and maintenance expenditure} & & \textbf{0} \\\hline
\end{tabularx}
\end{table}

\section{Total Estimated Project Cost}\label{cost-estimation-total-summary}

Table 4.8 consolidates actual project-specific cash expenditure. The total is zero because no project-specific purchase, paid deployment, external labour payment, or incremental service fee was incurred. COCOMO effort figures are kept separate because they describe estimated effort and schedule, not cash expenditure.

\begin{table}[H]
\centering\small
\caption{Total actual project cash expenditure}\label{tab:4-9}
\begin{tabularx}{\textwidth}{|>{\raggedright\arraybackslash}X|>{\raggedright\arraybackslash}p{1.55in}|>{\raggedleft\arraybackslash}p{1.25in}|}\hline
\textbf{Cost head} & \textbf{Source} & \textbf{Actual cost in \rupee} \\\hline
Human resources & Table 4.5 & 0 \\\hline
Hardware, software and infrastructure & Table 4.6 & 0 \\\hline
Testing, deployment and maintenance & Table 4.7 & 0 \\\hline
Project-specific service or API charges & No incremental amount paid from the project budget & 0 \\\hline
\textbf{Total actual project cash expenditure} & & \textbf{0} \\\hline
\end{tabularx}
\end{table}

\section{Assumptions and Cost Justification}\label{assumptions-and-cost-justification}

The zero total follows an actual-cash-expenditure basis. Only amounts paid specifically for the project are included; unpaid student effort, depreciation of equipment already owned, normal household or institute utilities, and possible future maintenance are not assigned hypothetical monetary values. Tools or services that were already available to the team and did not result in an additional project-specific payment are recorded as zero incremental expenditure. This accounting basis should not be confused with the COCOMO estimate, which estimates software effort and schedule rather than actual cash spending.


The treatment of each cost item follows from this basis. Student effort is recorded as zero because the work was unpaid and completed as part of the academic project. Existing equipment is zero because no new hardware was purchased. Software and developer tooling are zero because the project used open-source or already-available tools with no additional charge. Connectivity and utilities are zero because existing services were used and no project-specific amount was paid. Testing and deployment are zero because testing and the MVP ran locally, with no paid cloud environment. Future maintenance is excluded from the incurred cost because no post-project maintenance expense had been incurred at the time of reporting.

The final actual cash expenditure for VaultShield is therefore \rupee 0. This figure represents project-specific out-of-pocket expenditure and does not assign hypothetical market rates to unpaid student work or resources already available to the team.

\section{Risk Assessment and Mitigation}\label{risk-assessment-and-mitigation}

Risks were identified across the technical, security and project-management dimensions and reviewed at each iteration. Table 4.9 lists them with qualitative likelihood and impact ratings and the mitigation adopted or planned.

\begingroup\small

\Needspace{8\baselineskip}
\par\vspace{\TableCapGap}
\noindent\begin{minipage}{\textwidth}
\captionof{table}{Risk assessment and mitigation}\label{tab:4-11}
\end{minipage}
\par
\begin{longtable}[]{@{}|
  >{\raggedright\arraybackslash}p{(\columnwidth - 11\tabcolsep) * \real{0.2000}} |
  >{\raggedright\arraybackslash}p{(\columnwidth - 11\tabcolsep) * \real{0.2000}} |
  >{\raggedright\arraybackslash}p{(\columnwidth - 11\tabcolsep) * \real{0.2000}} |
  >{\raggedright\arraybackslash}p{(\columnwidth - 11\tabcolsep) * \real{0.2000}} |
  >{\raggedright\arraybackslash}p{(\columnwidth - 11\tabcolsep) * \real{0.2000}}|@{}}
\hline
\# & Risk & Likelihood & Impact & Mitigation \\
\hline
\endfirsthead
\hline
\# & Risk & Likelihood & Impact & Mitigation \\
\hline
\endhead
\endlastfoot
R1 & Container escape through a kernel or Docker vulnerability, widened in the default off mode where the sandbox runs as root with default capabilities & Low & Very high & Docker socket never mounted; only a read-only dummy credential file is mounted; risk documented as residual; non-root user, capability drop and gVisor or microVM isolation planned \\\hline
R2 & Malicious dependency executing install hooks inside the sandbox & Medium & High & Dependencies install only inside the disposable sandbox with dummy credentials; egress allowlist available; documented as residual risk \\\hline
R3 & Network exfiltration when egress is off & Medium & High & log-only and allowlist modes with a hardened proxy on an internal network; default-deny egress proposed as future work \\\hline
R4 & Exfiltration through an allowlisted host offering a general upload API & Low & High & Allowlist kept minimal per agent; the residual path is documented in the CLI README \\\hline
R5 & Real secrets reaching the sandbox or artifacts & Low & Very high & Shape-preserving dummy credentials; one-shot clone credential; token-hygiene scan with diff quarantine; scrubbed event payloads \\\hline
R6 & Agent tampering with Git state to hide changes from review & Low & High & Trusted capture in a separate helper from a pre-agent base read through the Docker API; unsafe Git configuration refused \\\hline
R7 & Approval gate bypassed through the unrestricted PATCH status=approved path by a token holder & Medium & Medium & Single-developer trust domain in the MVP; risk documented; restrict PATCH to gated transitions and add reviewer identities in future work \\\hline
R8 & Monitoring gaps & Medium & Medium & The trusted diff is the authoritative change record; watcher reconciliation pass; discrepancy detection planned \\\hline
R9 & Upstream agent CLI changes break an adapter, or allowlists drift & High & Medium & Adapter interface isolates changes; mock adapter keeps the pipeline testable; allowlist verification scheduled for Semester VIII \\\hline
R10 & Docker Desktop or platform behaviour changes & Medium & Medium & Docker-backed integration tests detect regressions \\\hline
R11 & MongoDB deployment without a replica set, breaking transactions & Medium & High & Backend documents the requirement; transaction tests fail loudly rather than skipping; compose stack points at Atlas \\\hline
R12 & Schedule slippage or member unavailability & Medium & Medium & Iterative delivery with a working pipeline after every iteration; overlapping module knowledge through shared testing \\\hline
R13 & Model API cost or rate limits during evaluation & Medium & Low & Deterministic mock agent for automated testing and benchmarking; real-agent runs limited to planned evaluation \\\hline\end{longtable}
\addtocounter{table}{-1}
\par\medskip

\endgroup

\clearpage

\chapter{Project Implementation}\label{project-implementation}

\section{Code Snippets}\label{code-snippets}

\subsection{Launching the Sandbox Container}\label{launching-the-sandbox-container}

The first snippet is the core of \texttt{Sandbox.\allowbreak{}create}, which creates the per-run container. Before any Docker call, it checks that the credential file exists, because Docker silently creates an empty directory at a missing bind-mount source; one historical run directory in the repository contains exactly such a directory, which illustrates why the check was added. The egress-related keyword arguments are passed only when they are set, so that in the default \texttt{off} mode the container is created exactly as before the egress feature existed. Any failure removes the partially created container and is reported as a single \texttt{SandboxError}.


\begin{figure}[H]
\centering
\SafeIncludeGraphics[width=\textwidth,height=0.8\textheight,keepaspectratio]{images/code/fig5-1-sandbox-create-teardown.png}
\par\medskip
\captionof{figure}{Sandbox container creation and teardown}\label{fig:5-1}
\end{figure}

As Figure 5.1 shows, the only host path given to the container is the read-only credential file; the repository is cloned inside the container afterwards. \texttt{teardown} uses \texttt{remove(v=True, force=True)}, which sends SIGKILL to the container and therefore to every process started in it, and it is called from a \texttt{finally} block in the \texttt{run} command so that it executes on every exit path.

\subsection{One-Shot Credential for Private Clones}\label{one-shot-credential-for-private-clones}

The second snippet shows how a private repository is cloned without giving the agent the repository-host token. The token, when available, is placed only in the \texttt{environment} of the single \texttt{git clone} exec, and the command line carries only the credential-helper configuration, which names environment variables rather than containing their values. \texttt{GIT\_\allowbreak{}TERMINAL\_\allowbreak{}PROMPT=0} makes an unauthenticated clone of a private repository fail immediately with Git's own message instead of hanging.


\begin{figure}[H]
\centering
\SafeIncludeGraphics[width=\textwidth,height=0.8\textheight,keepaspectratio]{images/code/fig5-2-one-shot-credential-clone.png}
\par\medskip
\captionof{figure}{Authenticated clone with a one-shot credential helper}\label{fig:5-2}
\end{figure}

The design in Figure 5.2 closes four leakage channels at once: the token is not in the remote URL, not in process arguments, not in the container environment, and not in error messages. Its lifetime is exactly the duration of \texttt{git clone}.

\subsection{Event Logging and Off-Limits Enforcement}\label{event-logging-and-off-limits-enforcement}

The third snippet shows the event record type, the off-limits path matcher and the append-only JSON Lines writer used by the monitor. Every record has exactly three fields, matching the backend's \texttt{RunEvent} model so that the uploader can forward lines without translation. The matcher applies both directory containment and glob matching against each spelling of the changed path, and it is deliberately generous.


\begin{figure}[H]
\centering
\SafeIncludeGraphics[width=\textwidth,height=0.8\textheight,keepaspectratio]{images/code/fig5-3-event-log-off-limits.png}
\par\medskip
\captionof{figure}{Event records, off-limits matching and the JSON Lines writer}\label{fig:5-3}
\end{figure}

In Figure 5.3, configuration order decides which entry is reported when a path violates several, which keeps the reported pattern deterministic. Because \texttt{json.\allowbreak{}dumps} escapes newlines and control characters inside a payload, each record always occupies exactly one line, which is what keeps a partially written log parseable. The monitor raises \texttt{OffLimitsViolation} only after this writer has run, so the violation record is always on disk before the run fails.

\subsection{Egress Allow/Deny Decision}\label{egress-allowdeny-decision}

The fourth snippet is the egress proxy's decision function. It is evaluated for every \texttt{CONNECT} request and every TLS connection identified by its Server Name Indication. The order of checks is significant: a host that is not on the allowlist is never resolved in \texttt{allowlist} mode, IP literals must match an IP entry exactly, and every resolved address must be globally routable, which prevents an allowlisted name from being pointed at an internal address.


\begin{figure}[H]
\centering
\SafeIncludeGraphics[width=\textwidth,height=0.8\textheight,keepaspectratio]{images/code/fig5-4-egress-decision.png}
\par\medskip
\captionof{figure}{Egress decision for one request target}\label{fig:5-4}
\end{figure}

Figure 5.4 also shows how \texttt{log-\allowbreak{}only} differs from \texttt{allowlist}: only the allowlist rule is relaxed; IP literals, port mismatches and private addresses are denied in both modes. Each decision is written to the proxy's log, which becomes \texttt{egress.\allowbreak{}jsonl} in the run directory and a \texttt{violation} event in the timeline.

\subsection{Branch Commit and Structured Diff}\label{branch-commit-and-structured-diff}

The fifth snippet shows the two Git operations at the heart of change capture, both of which run inside the trusted capture helper rather than the agent's sandbox. \texttt{commit\_\allowbreak{}run\_\allowbreak{}branch} refuses unsafe repository configurations before staging, then creates exactly one commit on the run branch. \texttt{build\_\allowbreak{}structured\_\allowbreak{}diff} obtains per-file counts with options that disable rename detection, external diff programs and text conversion filters, so that no configuration supplied by the agent can execute during diffing.


\begin{figure}[H]
\centering
\SafeIncludeGraphics[width=\textwidth,height=0.8\textheight,keepaspectratio]{images/code/fig5-5-run-branch-diff.png}
\par\medskip
\captionof{figure}{Committing the run branch and computing the diff}\label{fig:5-5}
\end{figure}

As shown in Figure 5.5, \texttt{-\allowbreak{}-\allowbreak{}allow-\allowbreak{}empty} guarantees that a run which changed nothing still produces a branch and an empty diff, because ``the agent changed nothing'' is a reviewable result rather than a missing artifact. \texttt{-\allowbreak{}-\allowbreak{}no-\allowbreak{}verify} prevents any hook in the repository from executing during the commit. If the per-file metadata alone would exceed the upload budget, the function fails before writing \texttt{diff.\allowbreak{}json}, because silently dropping entries would understate the change set.

\subsection{Approval State Machine}\label{approval-state-machine}

The sixth snippet is the inner function of \texttt{apply\_\allowbreak{}decision} in the backend, which implements the approve and reject endpoints. The status check and the status write are a single conditional update, and the \texttt{Approval} record is created in the same MongoDB transaction. The comment preserved in the source records the measurement that motivated the retry wrapper: before it was added, 54 of 72 losing calls in a concurrency test received a server error instead of the intended HTTP 409.


\begin{figure}[H]
\centering
\SafeIncludeGraphics[width=\textwidth,height=0.8\textheight,keepaspectratio]{images/code/fig5-6-approval-transaction.png}
\par\medskip
\captionof{figure}{Transactional approve/reject decision}\label{fig:5-6}
\end{figure}

Figure 5.6 shows why exactly one of two concurrent decisions can succeed: MongoDB applies an update to a single document atomically, so only one caller's filter on \texttt{status\_\allowbreak{}\_\allowbreak{}in=\{"awaiting\_\allowbreak{}review"\}} can match. The loser either sees zero documents updated or, after a transient write conflict, retries and finds the winner's \texttt{Approval}; both paths raise \texttt{StateTransitionError}, which the view renders as HTTP 409 naming the run's actual status. Because \texttt{Approval.\allowbreak{}run} is one-to-one, a second decision is also impossible at the schema level.

\subsection{Review Interface Decision Hook}\label{review-interface-decision-hook}

The seventh snippet is the dashboard function that records a reviewer's decision. It captures the run identifier at click time and ignores the response if the reviewer has since navigated to a different run, so that one run can never take on another run's status. When the backend refuses a decision, the status it reports is shown, so the reviewer sees the true state rather than a stale one.


\begin{figure}[H]
\centering
\SafeIncludeGraphics[width=\textwidth,height=0.8\textheight,keepaspectratio]{images/code/fig5-7-review-decision-hook.png}
\par\medskip
\captionof{figure}{Approve/reject handler in the run detail view}\label{fig:5-7}
\end{figure}

In Figure 5.7, \texttt{setPendingDecision} disables both buttons while a request is in flight, which prevents a reviewer from recording two decisions by double-clicking. \texttt{mergeDecidedRun} merges the decision response over the loaded run detail instead of replacing it, because the decision response carries only the run's flat fields; replacing the object would blank the four evidence tabs after a successful approval. The API client attaches \texttt{Authorization:\allowbreak{} Token <token>} to both calls.

\subsection{Submission Gate}\label{submission-gate}

The eighth snippet is the check that \texttt{vaultshield submit} performs before any side effect. It reads the run from the backend and refuses unless the status is exactly \texttt{approved}, naming the actual status, because ``not approved'' alone does not tell the developer what to do next.


\begin{figure}[H]
\centering
\SafeIncludeGraphics[width=\textwidth,height=0.8\textheight,keepaspectratio]{images/code/fig5-8-submit-approval-gate.png}
\par\medskip
\captionof{figure}{The approval gate in vaultshield submit}\label{fig:5-8}
\end{figure}

Figure 5.8 completes the approval path. A refused submission performs one \texttt{GET} and nothing else: no branch is pushed, no request is opened and no field is written. An approved submission then verifies that the checkout's \texttt{origin} matches the repository recorded in \texttt{branch.\allowbreak{}json}, imports the branch from the SHA-256-verified bundle with a ref-only fetch, pushes without \texttt{-\allowbreak{}-\allowbreak{}force}, opens the PR or MR, and records the resulting \texttt{pr\_\allowbreak{}url}.

\section{Steps to Access the System}\label{steps-to-access-the-system}

The prerequisites are listed in Table 5.1, followed by the installation and execution steps. The commands assume macOS or Linux with a POSIX shell.

\Needspace{8\baselineskip}
\par\vspace{\TableCapGap}
\noindent\begin{minipage}{\textwidth}
\captionof{table}{Software prerequisites}\label{tab:5-1}
\end{minipage}
\par
\begin{longtable}[]{@{}|
  >{\raggedright\arraybackslash}p{(\columnwidth - 7\tabcolsep) * \real{0.3333}} |
  >{\raggedright\arraybackslash}p{(\columnwidth - 7\tabcolsep) * \real{0.3333}} |
  >{\raggedright\arraybackslash}p{(\columnwidth - 7\tabcolsep) * \real{0.3333}}|@{}}
\hline
Tool & Version constraint & Purpose \\
\hline
\endfirsthead
\hline
Tool & Version constraint & Purpose \\
\hline
\endhead
\endlastfoot
Python & 3.11 or later & CLI and backend \\\hline
Docker & Running daemon & Sandbox, proxy and capture helper containers; Compose stack \\\hline
Git & Any recent version & Clone, branch, diff, bundle, push \\\hline
Node.js & 18 or later & Dashboard development and build \\\hline
MongoDB & Replica set or Atlas cluster & Backend datastore \\\hline
Playwright Chromium & One-time install & Screenshots of the running application \\\hline\end{longtable}
\addtocounter{table}{-1}
\par\medskip

\begin{enumerate}
\def\labelenumi{\arabic{enumi}.}
\item
  \textbf{Clone the repository.} git clone <repository-url>/VaultShield.git \&\& cd VaultShield
\item
  \textbf{Configure the backend environment.} Copy the template and fill in a MongoDB replica-set or Atlas connection string; never commit the resulting file. cp .env.example .env and set VAULTSHIELD\_\allowbreak{}MONGODB\_\allowbreak{}URL, optionally VAULTSHIELD\_\allowbreak{}MONGODB\_\allowbreak{}NAME, and, for anything shared, VAULTSHIELD\_\allowbreak{}SECRET\_\allowbreak{}KEY and VAULTSHIELD\_\allowbreak{}DEBUG=0.
\item
  \textbf{Start the backend and dashboard stack.} docker compose up --build This starts redis, web, channels, celery-worker, celery-beat and frontend. The API is served on http://localhost:8000 and runs migrations on start, and the dashboard is served on http://localhost:5173.
\item
  \textbf{Create a user account.} Self-service sign-up is not part of the MVP, so create a user through Django: docker compose exec web python manage.py createsuperuser
\item
  \textbf{Install the CLI.}

\begin{lstlisting}[language=bash]
cd cli
python3 -m venv .venv && source .venv/bin/activate
pip install -e ".[dev]"
playwright install chromium
docker build -t vaultshield-sandbox:local sandbox-image   # optional: Node + Antigravity image
\end{lstlisting}
\item
  \textbf{Authenticate the CLI}: vaultshield login --backend-url http://localhost:8000 --username <user>
\item
  \textbf{Create a project} in the dashboard with the New project button, and note its 24-character project identifier.
\item
  \textbf{Initialise the target repository's configuration.} In a checkout of the repository to be changed, run vaultshield init, which detects Django, Node, a Procfile or a Compose file and writes .vaultshield.yml with start\_\allowbreak{}command, port, ready\_path, routes and off\_\allowbreak{}limits\_\allowbreak{}paths. Optionally add project: <project-id> and an egress: block.
\item
  \textbf{Execute a sandboxed run.} The mock agent is the safest end-to-end test:

\begin{lstlisting}[language=bash]
vaultshield run --agent=mock --prompt="add a file" \
  --repo-url="$REPO" --base-image=vaultshield-sandbox:local \
  --project "$PROJECT_ID" [--egress=allowlist]
\end{lstlisting}

  A successful run prints the run identifier, ``Sandbox created'', ``Committed run \ldots{} to branch vaultshield/run-'', ``Sandbox torn down'' and ``Run \ldots{} complete'', uploads the artifacts and leaves the run in awaiting\_\allowbreak{}review.
\item
  \textbf{Inspect local artifacts}: open the newest run directory and view diff.json, events.jsonl and the captured screenshots.
\item
  \textbf{Review the run.} Open http://localhost:5173, sign in, select the project and the run, inspect the Diff, Timeline, Screenshots and Security tabs, and choose Approve or Reject.
\item
  \textbf{Submit or clean up.} For an approved run, from a checkout of the same repository: vaultshield submit <run-id>, which opens the PR or MR and prints its URL. For a rejected run: vaultshield teardown <run-id>.
\end{enumerate}

\section{System Implementation}\label{system-implementation}

This section describes each major interface of the implemented system. The dashboard screenshots listed below were captured by the repository's layout tool, which serves the built dashboard and answers every API request with fixture data. They therefore show the real interface rendering illustrative data. 

\subsection{Command-Line Interface and Run Directory}\label{command-line-interface-and-run-directory}

The CLI is the developer's primary interface. vaultshield --help lists the pipeline commands init, login, run, submit and teardown; the supporting commands ask, config, protect, status, exec, secure, unprotect, agents and hook; and a debugging command, \_sandbox-smoke, that creates and tears down a real sandbox against a repository URL. During vaultshield run, the terminal shows one line per pipeline milestone and, on any failure, a single-line error that names the stage and whether a sandbox was created. Each run writes its own run directory containing .env.sandbox, events.jsonl, diff.json, branch.bundle, branch.json, the screenshots and, when egress is enabled, egress.jsonl.


\Needspace{22\baselineskip}
\noindent\begin{minipage}{\textwidth}
\centering
\fbox{\SafeIncludeGraphics[width=\textwidth,height=0.82\textheight,keepaspectratio]{images/screenshots/terminal-run.png}}
\captionof{figure}{Terminal output of a successful vaultshield run with the mock agent}\label{fig:5-9}
\end{minipage}
\par\medskip

Figure 5.9 shows these milestone lines for one complete run, from registration through screenshot capture, the run-branch commit, sandbox teardown and upload. Figure 5.10 is a real artifact: the screenshot captured for route / in run 20261009-092431-6891a631b5f4a313, in which the target application was a static file server started inside the sandbox.


\Needspace{22\baselineskip}
\noindent\begin{minipage}{\textwidth}
\centering
\fbox{\SafeIncludeGraphics[width=0.9\textwidth,keepaspectratio]{images/screenshots/root.png}}
\captionof{figure}{Screenshot artifact captured inside the sandbox}\label{fig:5-10}
\end{minipage}
\par\medskip

The captured page in Figure 5.10 is the directory listing of the cloned repository served by python -m http.server, and it includes vaultshield\_\allowbreak{}mock\_\allowbreak{}agent.txt, the file the mock agent created. It therefore demonstrates visually that the screenshot was taken inside the sandbox, after the agent ran, against the sandbox's clone rather than the developer's checkout.

\subsection{Landing Page}\label{landing-page}

The landing page introduces the product and its supporting Protect module to new users. It shows the headline ``Let AI agents read your code, not your secrets.'', a one-line installer command with a copy button, navigation to ``How it works'' and the FAQ, a theme switch and links to sign in. The installer script and the wheel it downloads are served from the dashboard's own origin.


\Needspace{22\baselineskip}
\noindent\begin{minipage}{\textwidth}
\centering
\fbox{\SafeIncludeGraphics[width=0.9\textwidth,keepaspectratio]{images/screenshots/landing-1280-light.png}}
\captionof{figure}{Landing page}\label{fig:5-11}
\end{minipage}
\par\medskip

\subsection{Sign-In Page}\label{sign-in-page}

The sign-in page uses a two-column layout: the left column restates the product's principle with a short terminal illustration of a protected .env value, and the right column contains the username and password form. Submitting the form calls POST /api/auth/login/, stores the returned token in localStorage, and opens the dashboard. Any later HTTP 401 response returns the user to this page.


\Needspace{22\baselineskip}
\noindent\begin{minipage}{\textwidth}
\centering
\fbox{\SafeIncludeGraphics[width=0.9\textwidth,keepaspectratio]{images/screenshots/login-1280-light.png}}
\captionof{figure}{Sign-in page}\label{fig:5-12}
\end{minipage}
\par\medskip

\subsection{Projects Dashboard}\label{projects-dashboard}

After sign-in, the projects overview shows one card per project with its name, repository URL and creation date, alongside a persistent sidebar listing the same projects, a New project button, a sign-out link and the theme switch. Selecting a project opens its run history. The New project button opens a modal dialog that collects the project name and repository URL and calls POST /api/projects/; a separate confirmation dialog guards Delete project, which removes the project together with its runs and artifacts.


\Needspace{22\baselineskip}
\noindent\begin{minipage}{\textwidth}
\centering
\fbox{\SafeIncludeGraphics[width=0.9\textwidth,keepaspectratio]{images/screenshots/dashboard-1280-light.png}}
\captionof{figure}{Projects dashboard}\label{fig:5-13}
\end{minipage}
\par\medskip


\Needspace{22\baselineskip}
\noindent\begin{minipage}{\textwidth}
\centering
\fbox{\SafeIncludeGraphics[width=0.9\textwidth,keepaspectratio]{images/screenshots/new-project-modal-1280-light.png}}
\captionof{figure}{New project modal dialog}\label{fig:5-14}
\end{minipage}
\par\medskip

\subsection{Run History}\label{run-history}

The run history page lists a project's runs newest first in a full-width table with the run identifier, agent, status badge and relative creation time. Filter chips select one status or all, a search box filters by run identifier, and a Load more control follows the backend's pagination. Clicking a run identifier opens the run detail page, with the active tab reflected in the URL so that a specific view can be shared.


\Needspace{22\baselineskip}
\noindent\begin{minipage}{\textwidth}
\centering
\fbox{\SafeIncludeGraphics[width=0.9\textwidth,keepaspectratio]{images/screenshots/runs-1280-light.png}}
\captionof{figure}{Run history with status filters}\label{fig:5-15}
\end{minipage}
\par\medskip

\subsection{Run Detail: Diff View}\label{run-detail-diff-view}

The run detail page shows the run identifier, agent, status and date in its header, a decision banner with Approve and Reject buttons when the run is in awaiting\_\allowbreak{}review, and four tabs. The Diff tab lists every changed file from diff.json with its addition and deletion counts, which are labelled for assistive technology because +12 and -3 alone are not meaningful to a screen reader. In the current release the tab does not render the patch text itself; the patches are stored with the diff and consumed by the security scan, and a reviewer who needs line-level content reads the patch fields in the run directory's diff.json or fetches the run branch from branch.bundle with standard Git commands. Rendering patches inline is identified as future work in Chapter 9.


\Needspace{22\baselineskip}
\noindent\begin{minipage}{\textwidth}
\centering
\fbox{\SafeIncludeGraphics[width=0.9\textwidth,keepaspectratio]{images/screenshots/detail-diff-1280-light.png}}
\captionof{figure}{Run detail – Diff tab}\label{fig:5-16}
\end{minipage}
\par\medskip

\subsection{Run Detail: Timeline View}\label{run-detail-timeline-logs-view}

The Timeline tab renders the uploaded Event\_Log in sequence order. Each entry shows its timestamp, an event-type badge and the payload, such as a command and its result, a modified file path, or an egress denial. This is the reviewer's view of what the agent did, as opposed to what it changed, and it is the primary aid to understanding a change in the sense identified by Bacchelli and Bird \cite{ref14}.


\Needspace{22\baselineskip}
\noindent\begin{minipage}{\textwidth}
\centering
\fbox{\SafeIncludeGraphics[width=0.9\textwidth,keepaspectratio]{images/screenshots/detail-timeline-1280-light.png}}
\captionof{figure}{Run detail – Timeline tab}\label{fig:5-17}
\end{minipage}
\par\medskip

\subsection{Run Detail: Screenshots View}\label{run-detail-screenshots-view}

The Screenshots tab displays one image per configured route, served from the backend's media storage. When a run configured no routes, or the application never became ready, the tab states that no screenshots were captured, so that an empty tab is never mistaken for a loading failure.


\Needspace{22\baselineskip}
\noindent\begin{minipage}{\textwidth}
\centering
\fbox{\SafeIncludeGraphics[width=0.9\textwidth,keepaspectratio]{images/screenshots/detail-screenshots-1280-light.png}}
\captionof{figure}{Run detail – Screenshots tab}\label{fig:5-18}
\end{minipage}
\par\medskip

\subsection{Run Detail: Security Findings View}\label{run-detail-security-findings-view}

The Security tab lists the findings produced by the Celery security scan of the diff's added lines, each with its rule identifier, file path and line number, and a description. The fixture shown uses the label hardcoded-secret; the rule identifier stored by the implemented scanner for this category is secret-literal. Findings inform the reviewer but do not block approval, consistent with the principle that the decision belongs to the human.


\Needspace{22\baselineskip}
\noindent\begin{minipage}{\textwidth}
\centering
\fbox{\SafeIncludeGraphics[width=0.9\textwidth,keepaspectratio]{images/screenshots/detail-security-1280-light.png}}
\captionof{figure}{Run detail – Security tab}\label{fig:5-19}
\end{minipage}
\par\medskip

\subsection{Approval Decision and Submission}\label{approval-decision-and-submission}

When the reviewer approves a run, the banner's buttons are replaced by the status copy ``approved, awaiting submission'' and the exact command vaultshield submit <run-id>, because the backend never opens a PR on its own; that action requires the developer's local repository credentials. For a rejected run the buttons disappear permanently, since the run can never be decided again. A run uploaded by vaultshield protect sync shows ``Protection log: nothing to review'' in place of the buttons.


\Needspace{22\baselineskip}
\noindent\begin{minipage}{\textwidth}
\centering
\fbox{\SafeIncludeGraphics[width=0.9\textwidth,keepaspectratio]{images/screenshots/detail-approved-1280-light.png}}
\captionof{figure}{Run detail after approval, showing the submit command}\label{fig:5-20}
\end{minipage}
\par\medskip

Figure 5.20 shows a real run immediately after the Approve button was pressed: the status badge reads ``approved, awaiting submission'', the banner gives the exact submit command, and the Diff tab still lists the one file the agent created.

\subsection{Implementation Summary}\label{implementation-summary}

The implementation combines a host-side command-line pipeline with a review dashboard. The CLI creates the sandbox, monitors the agent, captures the change on a dedicated branch and uploads the evidence, while the dashboard presents that evidence through the Diff, Timeline, Screenshots and Security tabs and records the reviewer's decision. Keeping execution in the CLI and review in the dashboard means that producing a change and accepting it remain separate steps, and only an approved run can be submitted as a pull or merge request.

\section{Timeline Sem VII}\label{project-timeline}

The project timeline presents the major milestones in the development of VaultShield during the seventh semester. The Gantt chart in Figure~5.21 gives a graphical overview of the activities and their scheduled duration across thirteen weeks, from group formation on 10 July 2026 to the external presentation and viva voce on 2 November 2026.

The activities are grouped into four phases. Project Conception and Initiation covers group formation and topic finalisation, discussion of the topic with a paper prototype, identification of the project's functionalities, the literature review and the design of the graphical user interface, ending with Presentation I. Project Design and Implementation covers policy logic development and core development, ending with Presentation II. Project Evaluation and Documentation covers incorporating the suggestions from Presentation II and submitting a draft to the project guide for proofreading. Final Submission and Enhancement covers the draft review, the final black book submission, and the external presentation and viva voce.

For each activity, the chart records the task owners, start and due dates, duration in weeks and percentage of completion. All four members are listed as owners of every activity, and every Semester VII activity is marked 100\% complete. The Semester VIII activities are listed in Table 4.1.

\clearpage
\begin{figure}[p]
\centering
\rotatebox{90}{%
  \begin{minipage}{0.92\textheight}
    \centering
    \SafeIncludeGraphics[width=\linewidth,height=0.95\textwidth,keepaspectratio]{images/diagrams/fig4-1-gantt-chart.png}
    \par\medskip
    \captionof{figure}{Gantt chart of the VaultShield major project}\label{fig:5-21}
  \end{minipage}%
}
\end{figure}
\clearpage

\clearpage

\chapter{Testing}\label{software-testing}

\section{Software Testing}\label{software-testing-1}

The testing strategy was shaped by the nature of the system. VaultShield's value lies almost entirely in negative guarantees: the agent must not reach the host, real secrets must not enter the sandbox, a token must not appear in any artifact, an unapproved run must not become a pull request. Negative guarantees are poorly served by a handful of example-based tests, because an example demonstrates only that one input behaves correctly. The team therefore adopted four principles. First, every requirement's acceptance criteria and every numbered correctness property in the design document was traced to at least one test, with the requirement or property identifier cited in the test's docstring. Second, invariants were expressed as property-based tests using Hypothesis and fast-check, which generate hundreds of inputs per property and shrink any counterexample to a minimal case.

Third, isolation claims were tested against real Docker rather than mocks. Docker-gated tests create real containers from the real sandbox image, clone a real repository served over git:// from the host, run a probing agent inside the container that actively attempts to read the host .env, list the host home directory and find the Docker socket, and then assert on what the probe could and could not see. These tests skip automatically when no Docker daemon is available, so the suite remains runnable anywhere, but in the campaign reported here Docker was available and no test was skipped. Fourth, the backend was tested against a real MongoDB replica set rather than an in-memory substitute, because the approval gate's correctness depends on genuine transaction and single-document atomicity semantics; a standalone mongod deliberately fails the transaction tests rather than skipping them. For this campaign a throwaway MongoDB 7 single-node replica set was started in a local container and removed afterwards, so no shared database was touched.

External services that cannot be exercised safely in tests were replaced with in-memory fakes at clearly defined seams: the macOS Keychain, the model API used by the orchestrator, and the backend during CLI pipeline tests. Test fixtures also redirect HOME to an empty temporary directory so that no test can read or write the developer's real home directory, agent configuration or Keychain.

\section{Types of Testing}\label{types-of-testing}

Testing was organised into two types, functional testing and system testing. Table 6.1 lists the five test levels used within them, and each type is described below.

\Needspace{8\baselineskip}
\par\vspace{\TableCapGap}
\noindent\begin{minipage}{\textwidth}
\captionof{table}{Types of testing performed}\label{tab:6-1}
\end{minipage}
\par
\begin{longtable}[]{@{}|
  >{\raggedright\arraybackslash}p{(\columnwidth - 9\tabcolsep) * \real{0.2500}} |
  >{\raggedright\arraybackslash}p{(\columnwidth - 9\tabcolsep) * \real{0.2500}} |
  >{\raggedright\arraybackslash}p{(\columnwidth - 9\tabcolsep) * \real{0.2500}} |
  >{\raggedright\arraybackslash}p{(\columnwidth - 9\tabcolsep) * \real{0.2500}}|@{}}
\hline
Type & Scope & Tools & Representative tests \\
\hline
\endfirsthead
\hline
Type & Scope & Tools & Representative tests \\
\hline
\endhead
\endlastfoot
Unit & Individual functions and classes: configuration parsing, run-ID generation, credential shaping, off-limits matching, diff parsing, egress decisions, serializers, React components & pytest, Hypothesis, Jest, Testing Library, fast-check & test\_\allowbreak{}config, test\_\allowbreak{}run\_\allowbreak{}directory, test\_\allowbreak{}credentials, test\_\allowbreak{}egress, api.test.js, TabPanel.test.jsx \\\hline
Integration & Interaction between modules and with real infrastructure: CLI with Docker, backend with MongoDB, upload with the API & pytest with real Docker and a MongoDB replica set & test\_\allowbreak{}sandbox\_\allowbreak{}execution, test\_\allowbreak{}egress\_\allowbreak{}docker, test\_\allowbreak{}trusted\_\allowbreak{}capture, test\_\allowbreak{}mongodb.py \\\hline
System & The complete vaultshield run command from configuration to teardown, and the submit and teardown flows & Click CliRunner, Docker, local git daemon; timing harness & test\_\allowbreak{}agent\_\allowbreak{}executes\_\allowbreak{}inside\_\allowbreak{}the\_\allowbreak{}sandbox\_\allowbreak{}and\_\allowbreak{}cannot\_\allowbreak{}reach\_\allowbreak{}the\_\allowbreak{}host, test\_\allowbreak{}run\_\allowbreak{}with\_\allowbreak{}egress\_\allowbreak{}allowlist\_\allowbreak{}end\_\allowbreak{}to\_\allowbreak{}end, test\_\allowbreak{}submits\_\allowbreak{}an\_\allowbreak{}approved\_\allowbreak{}run\_\allowbreak{}and\_\allowbreak{}returns\_\allowbreak{}the\_\allowbreak{}pull\_\allowbreak{}request\_\allowbreak{}url \\\hline
Security-focused & Isolation, secret suppression, network blocking, tamper resistance, token hygiene, authentication, approval atomicity & Probing agents in real containers; property tests; concurrency tests with multiple workers & test\_\allowbreak{}dns\_\allowbreak{}exfil\_\allowbreak{}is\_\allowbreak{}blocked, test\_\allowbreak{}a\_\allowbreak{}tampered\_\allowbreak{}in\_\allowbreak{}container\_\allowbreak{}git\_\allowbreak{}cannot\_\allowbreak{}hide\_\allowbreak{}a\_\allowbreak{}file\_\allowbreak{}from\_\allowbreak{}the\_\allowbreak{}branch\_\allowbreak{}or\_\allowbreak{}diff, test\_\allowbreak{}token\_\allowbreak{}hygiene\_\allowbreak{}vectors\_\allowbreak{}in\_\allowbreak{}a\_\allowbreak{}real\_\allowbreak{}container, test\_\allowbreak{}concurrent\_\allowbreak{}decisions\_\allowbreak{}on\_\allowbreak{}one\_\allowbreak{}run\_\allowbreak{}yield\_\allowbreak{}exactly\_\allowbreak{}one\_\allowbreak{}approval \\\hline
Usability and accessibility & Contrast ratios, responsive layout, keyboard navigation, status copy & Jest contrast test; Playwright layout tool over 9 pages × 7 widths & contrast.test.js, the layout report, RunDetailView.test.jsx \\\hline\end{longtable}
\addtocounter{table}{-1}
\par\medskip

\subsection{Functional Testing}\label{functional-testing}

Functional testing verified that each function, module and security control behaves as specified. It was carried out at three levels: unit, integration and security-focused testing.

Unit testing covered pure logic in isolation. Examples include the property that a generated run identifier always satisfies the run-ID validator and is safe as a file name and Git ref component; the property that a dummy credential never contains the real value as a substring; the parsing of git diff --numstat -z output, including binary files and non-UTF-8 file names; and the egress decision table for IP literals, port mismatches, wildcard entries and private addresses. On the frontend, unit tests verify that the API client attaches the Token header to every call except login, that status copy for an approved run is exactly ``approved, awaiting submission'', and that the tab list honours arrow-key navigation.

Integration testing exercised boundaries that mocks would misrepresent. The CLI's Docker-backed tests verify, among other things, that an exec's environment does not persist into later execs, that the credential helper actually feeds Git inside the container, that the sandbox image provides every binary the pipeline executes and carries no credentials, and that a nested project\_\allowbreak{}dir really becomes the agent's working directory. The backend's MongoDB integration suite verifies REST round trips, foreign-key behaviour, atomic sequence allocation under concurrency and transactional rollback.

Security-focused testing was the largest category by intent. Beyond the isolation probe, it includes blocking and recording of a non-allowlisted connection, blocking of DNS-based exfiltration from the internal network, reachability of allowlisted hosts and of the application port through the proxy, five token-hygiene vectors in a real container, refusal to follow symlinks out of the credential root, and a tampered in-container Git that attempts to hide a file from the branch and diff. On the backend, it includes rejection of unauthenticated requests with no side effects and concurrent approve/reject requests resolving to exactly one decision.

\subsection{System Testing}\label{system-testing}

System testing drove the real vaultshield run command end to end. The principal system test runs the complete pipeline with a real container, a real clone, a real in-container agent process, real trusted capture and real teardown, substituting only the adapter so that a probe can report from inside the sandbox. The same approach underlies the timing harness used in Chapter 7, which executed 40 additional complete runs.

Usability and accessibility testing combined automated checks with layout inspection. A Jest test computes WCAG contrast ratios for every foreground/background token pair in both themes; the layout tool renders nine pages at seven widths from 390 px to 2560 px, records horizontal overflow and container bounds, and saves screenshots for review. A formal usability study with external reviewers was not conducted in Semester VII and is planned for Semester VIII.

\section{Testing Results and Findings}\label{testing-results-and-findings}

Table 6.2 summarises the results of the three automated suites.

\begingroup\scriptsize

\Needspace{8\baselineskip}
\par\vspace{\TableCapGap}
\noindent\begin{minipage}{\textwidth}
\captionof{table}{Automated test suite results}\label{tab:6-2}
\end{minipage}
\par
\begin{longtable}[]{@{}|
  >{\raggedright\arraybackslash}p{(\columnwidth - 17\tabcolsep) * \real{0.0968}} |
  >{\raggedleft\arraybackslash}p{(\columnwidth - 17\tabcolsep) * \real{0.1290}} |
  >{\raggedleft\arraybackslash}p{(\columnwidth - 17\tabcolsep) * \real{0.1290}} |
  >{\raggedleft\arraybackslash}p{(\columnwidth - 17\tabcolsep) * \real{0.1290}} |
  >{\raggedleft\arraybackslash}p{(\columnwidth - 17\tabcolsep) * \real{0.1290}} |
  >{\raggedleft\arraybackslash}p{(\columnwidth - 17\tabcolsep) * \real{0.1290}} |
  >{\raggedleft\arraybackslash}p{(\columnwidth - 17\tabcolsep) * \real{0.1290}} |
  >{\raggedleft\arraybackslash}p{(\columnwidth - 17\tabcolsep) * \real{0.1290}}|@{}}
\hline
Suite & Test files & Tests & Passed & Failed & Skipped & Pass rate & Wall-clock time \\
\hline
\endfirsthead
\hline
Suite & Test files & Tests & Passed & Failed & Skipped & Pass rate & Wall-clock time \\
\hline
\endhead
\endlastfoot
CLI & 52 & 2,080 & 2,078 & 2 & 0 & 99.90\% & 731.4 s* \\\hline
Backend & 6 & 341 & 341 & 0 & 0 & 100.00\% & 21.3 s \\\hline
Frontend & 21 & 283 & 283 & 0 & 0 & 100.00\% & 2.4 s \\\hline
\textbf{Total} & \textbf{79} & \textbf{2,704} & \textbf{2,702} & \textbf{2} & \textbf{0} & \textbf{99.93\%} & -- \\\hline\end{longtable}
\addtocounter{table}{-1}
\par\medskip

\endgroup

*600.3 s of the CLI suite's time was consumed by a single failing test that waited for its fake sign-in program's ten-minute sleep to end; the remaining 2,079 tests completed in about 131 s.

Table 6.3 breaks the CLI suite down by functional area.

\begingroup\small

\Needspace{8\baselineskip}
\par\vspace{\TableCapGap}
\noindent\begin{minipage}{\textwidth}
\captionof{table}{CLI test results by functional area}\label{tab:6-3}
\end{minipage}
\par
\begin{longtable}[]{@{}|
  >{\raggedright\arraybackslash}p{(\columnwidth - 11\tabcolsep) * \real{0.1667}} |
  >{\raggedright\arraybackslash}p{(\columnwidth - 11\tabcolsep) * \real{0.1667}} |
  >{\raggedleft\arraybackslash}p{(\columnwidth - 11\tabcolsep) * \real{0.2222}} |
  >{\raggedleft\arraybackslash}p{(\columnwidth - 11\tabcolsep) * \real{0.2222}} |
  >{\raggedleft\arraybackslash}p{(\columnwidth - 11\tabcolsep) * \real{0.2222}}|@{}}
\hline
Area & Test modules & Tests & Passed & Failed \\
\hline
\endfirsthead
\hline
Area & Test modules & Tests & Passed & Failed \\
\hline
\endhead
\endlastfoot
Pipeline, configuration and run directory & test\_cli, test\_\allowbreak{}config, test\_\allowbreak{}project\_\allowbreak{}dir, test\_\allowbreak{}run\_\allowbreak{}directory, test\_\allowbreak{}install\_\allowbreak{}sh & 293 & 293 & 0 \\\hline
Sandbox, credentials and isolation & test\_\allowbreak{}sandbox\_\allowbreak{}execution, test\_\allowbreak{}sandbox\_\allowbreak{}clone\_\allowbreak{}credentials, test\_\allowbreak{}credentials, test\_\allowbreak{}interactive\_\allowbreak{}auth, test\_\allowbreak{}antigravity & 259 & 257 & 2 \\\hline
Egress control & test\_\allowbreak{}egress, test\_\allowbreak{}egress\_\allowbreak{}docker & 92 & 92 & 0 \\\hline
Monitoring and evidence & test\_\allowbreak{}monitor, test\_\allowbreak{}browser\_\allowbreak{}capture & 58 & 58 & 0 \\\hline
Capture and token hygiene & test\_\allowbreak{}git\_\allowbreak{}capture, test\_\allowbreak{}trusted\_\allowbreak{}capture, test\_\allowbreak{}token\_\allowbreak{}hygiene & 170 & 170 & 0 \\\hline
Upload, submit, teardown and providers & test\_\allowbreak{}uploader, test\_\allowbreak{}submit, test\_\allowbreak{}teardown, test\_\allowbreak{}providers & 127 & 127 & 0 \\\hline
Orchestrate & 9 test\_\allowbreak{}orchestrate\_\allowbreak{}* modules & 377 & 377 & 0 \\\hline
Protect & 22 test\_\allowbreak{}protect\_\allowbreak{}* modules & 704 & 704 & 0 \\\hline
\textbf{Total} & \textbf{52} & \textbf{2,080} & \textbf{2,078} & \textbf{2} \\\hline\end{longtable}
\addtocounter{table}{-1}
\par\medskip

\endgroup

The backend's 341 tests were distributed as follows: 154 MongoDB integration tests, 72 run API and state-machine tests, 43 security-scan tests, 37 scheduled-task tests, 22 settings and URL tests, and 13 project API tests. Of the 39 Docker-backed CLI tests, 37 passed. The accessibility test reported a lowest text contrast of 4.77:1 in the light theme and 5.80:1 in the dark theme, both above the 4.5:1 threshold, and a lowest non-text UI contrast of 3.37:1 and 3.30:1 respectively, both above the 3:1 threshold for interface components. The layout report recorded 63 page-width combinations with zero horizontal overflow and zero layout problems.

Table 6.4 lists the issues found during development and during this campaign. Issues I-01 to I-06 and I-10 were found and fixed during development and are documented in the source code or README; issues I-07 to I-09 and I-11 were found or confirmed during the campaign and remain open.

\begingroup\small

\Needspace{8\baselineskip}
\par\vspace{\TableCapGap}
\noindent\begin{minipage}{\textwidth}
\captionof{table}{Issues found and their resolution}\label{tab:6-4}
\end{minipage}
\par
\begin{longtable}[]{@{}|
  >{\raggedright\arraybackslash}p{(\columnwidth - 11\tabcolsep) * \real{0.2000}} |
  >{\raggedright\arraybackslash}p{(\columnwidth - 11\tabcolsep) * \real{0.2000}} |
  >{\raggedright\arraybackslash}p{(\columnwidth - 11\tabcolsep) * \real{0.2000}} |
  >{\raggedright\arraybackslash}p{(\columnwidth - 11\tabcolsep) * \real{0.2000}} |
  >{\raggedright\arraybackslash}p{(\columnwidth - 11\tabcolsep) * \real{0.2000}}|@{}}
\hline
ID & Issue & How it was found & Resolution & Status \\
\hline
\endfirsthead
\hline
ID & Issue & How it was found & Resolution & Status \\
\hline
\endhead
\endlastfoot
I-01 & Concurrent approve/reject: most losing requests received HTTP 500 instead of 409; measured at 54 of 72 losing calls & Concurrency test after the MongoDB migration & Retry on TransientTransactionError; refusal path returns 409 with the actual status & Fixed \\\hline
I-02 & Dummy connection strings preserved the real host name, leaking a real MongoDB Atlas cluster host into .env.sandbox & Observed in a .env.sandbox generated from a real .env & Host replaced unless provably local; property tests added & Fixed \\\hline
I-03 & A missing credential-file path was silently turned into an empty directory by Docker's bind mount & Failure in a real run & Path.is\_\allowbreak{}file() check before any Docker call & Fixed \\\hline
I-04 & Docker VM listeners were reachable from a plain internal bridge & Port measurement from a container on a plain internal bridge & Internal bridge created without a gateway address & Fixed \\\hline
I-05 & A private clone without a credential hung or failed with an opaque exit code & Live failure reproduced in a Docker test & GIT\_\allowbreak{}TERMINAL\_\allowbreak{}PROMPT=0 and an actionable error naming the missing credential & Fixed \\\hline
I-06 & Antigravity sign-in failed under the allowlist without www.googleapis.com, but allowing it also permits that host's upload APIs & Real log-only and allowlist probes & Host added with the residual exfiltration risk documented & Mitigated \\\hline
I-07 & The in-sandbox sign-in read deadline did not end a sign-in that prompts and then hangs; the test waited 600.3 s and the auth phase returned success & test\_\allowbreak{}a\_\allowbreak{}hanging\_\allowbreak{}signin\_\allowbreak{}hits\_\allowbreak{}the\_\allowbreak{}read\_\allowbreak{}timeout\_\allowbreak{}instead\_\allowbreak{}of\_\allowbreak{}blocking\_\allowbreak{}forever & Under investigation; suspected interaction between the interactive exec socket and Docker Desktop 29.2.1 & Open \\\hline
I-08 & Interactive exec test could not observe its stdin sentinel because the sentinel was redacted from the container's echoed output & test\_\allowbreak{}exec\_\allowbreak{}interactive\_\allowbreak{}really\_\allowbreak{}delivers\_\allowbreak{}stdin\_\allowbreak{}to\_\allowbreak{}a\_\allowbreak{}container\_\allowbreak{}process & Redaction erred on the safe side; test expectation and scrubbing behaviour to be reconciled & Open \\\hline
I-09 & The in-container file watcher missed the agent's file creation in some runs, although the trusted diff recorded the file every time & Benchmark and inspection of historical run directories & Diff remains the authoritative record; cross-checking the event log against the diff is planned & Open \\\hline
I-10 & Frontend Docker image failed to build because of a stale lock file & docker compose build frontend & Lock file regenerated & Fixed \\\hline
I-11 & PATCH /api/runs/:id/ accepts status=approved without an Approval record, so any token holder could bypass the dashboard decision & Code review & Documented as a known limitation of the single-developer MVP; gating all status transitions is planned & Open \\\hline\end{longtable}
\addtocounter{table}{-1}
\par\medskip

\endgroup

\section{Sample Test Cases}\label{sample-test-cases}

Table 6.5 presents representative test cases. Each corresponds to an automated test or to the benchmark harness of Chapter 7, and the ``actual result'' column reports the observed outcome.

\begingroup\small

\Needspace{8\baselineskip}
\par\vspace{\TableCapGap}
\noindent\begin{minipage}{\textwidth}
\captionof{table}{Sample test cases}\label{tab:6-5}
\end{minipage}
\par
\begin{longtable}[]{@{}|
  >{\raggedright\arraybackslash}p{(\columnwidth - 13\tabcolsep) * \real{0.1667}} |
  >{\raggedright\arraybackslash}p{(\columnwidth - 13\tabcolsep) * \real{0.1667}} |
  >{\raggedright\arraybackslash}p{(\columnwidth - 13\tabcolsep) * \real{0.1667}} |
  >{\raggedright\arraybackslash}p{(\columnwidth - 13\tabcolsep) * \real{0.1667}} |
  >{\raggedright\arraybackslash}p{(\columnwidth - 13\tabcolsep) * \real{0.1667}} |
  >{\raggedright\arraybackslash}p{(\columnwidth - 13\tabcolsep) * \real{0.1667}}|@{}}
\hline
ID & Description & Input & Expected output & Actual result & Status \\
\hline
\endfirsthead
\hline
ID & Description & Input & Expected output & Actual result & Status \\
\hline
\endhead
\endlastfoot
TC-01 & Sandbox isolation: no host writes and no host access & Full vaultshield run with a probing agent that edits files, reads .env from its working directory, lists the host home directory and looks for the Docker socket & Agent working directory is the sandbox clone; host .env unreadable and secret not seen; home not listable; socket absent; host files, .env and HEAD unchanged; no run branch in host repo; container removed & All assertions held; in the benchmark, the host checkout was untouched in 40 of 40 runs & Pass \\\hline
TC-02 & Secret suppression: dummy values never contain the real value & Hypothesis-generated URLs, keys, connection strings and generic values & Same keys in the same order; no dummy contains its real value; userinfo replaced; non-local hosts replaced; local hosts and ports preserved & 135 of 135 tests passed & Pass \\\hline
TC-03 & Repository token confined to the clone & Exec with an environment variable, followed by a second exec; authenticated clone through the credential helper & Variable absent from the second exec; Git receives the credential through the helper with the token absent from argv & Both passed & Pass \\\hline
TC-04 & Network blocking in allowlist mode & From inside the isolated sandbox: CONNECT to example.org:443, to the IP literal 1.1.1.1:443 and to the host at host.\allowbreak{}docker.\allowbreak{}internal:8000; a direct DNS lookup of a unique label, then a CONNECT to it & Every CONNECT answered HTTP/1.1 403 by the proxy; the direct DNS lookup fails because the internal network has no resolver; proxy log records deny / not\_\allowbreak{}in\_\allowbreak{}allowlist & Both passed; an allowlisted host remained reachable & Pass \\\hline
TC-05 & Off-limits enforcement & Agent writes to protected.txt, configured in off\_\allowbreak{}limits\_\allowbreak{}paths & violation event written to the Event\_Log, then the run fails with OffLimitsViolation; sandbox removed & Passed & Pass \\\hline
TC-06 & Diff correctness & Generated file edits with known added and removed lines; 40 mock-agent runs & Counts in diff.json equal the + and - lines of each patch; the agent's file appears in every run's diff & Passed; agent's file present in 40 of 40 benchmark diffs & Pass \\\hline
TC-07 & Tamper resistance & The sandbox's own Git is tampered with so that it would hide a file the agent created & The file still appears on the run branch with a regular-file mode and no gitlink entries, because the branch is built by the helper from the trusted base and a filtered archive & Both variants passed & Pass \\\hline
TC-08 & Token hygiene & Agent copies its credential into the repo, redirects core.worktree, plants a symlink, writes a Git-ignored copy, or commits a tracked symlink & Run fails with TokenHygieneViolation; diff renamed diff.\allowbreak{}quarantined.\allowbreak{}json; no bundle written; message never contains the secret & All five vectors passed & Pass \\\hline
TC-09 & Approval gating: unapproved changes never reach a PR & vaultshield submit for runs in awaiting\_\allowbreak{}review, rejected, running and approved & Non-approved: SubmitError naming the status, no push, no PR, no write; approved: push, PR opened, pr\_url recorded & Passed & Pass \\\hline
TC-10 & Concurrent decisions & 6 simultaneous approve/reject calls on one run; 4 further trials of 5 simultaneous approvals & Exactly one winner per trial; every loser receives a clean refusal; exactly one Approval row & Passed against the MongoDB replica set & Pass \\\hline
TC-11 & Authentication required & API requests without an Authorization header & HTTP 401; no document created or modified & Passed & Pass \\\hline
TC-12 & Monitoring completeness: every agent write appears in the timeline & 40 mock-agent runs, each creating one file; unit test with interleaved producers & Every run's Event\_Log contains a file\_\allowbreak{}change event for the agent's file; merged log is timestamp-ordered & Ordering test passed; file\_\allowbreak{}change present in 32 of 40 runs; the diff recorded the change in 40 of 40 & Partial \\\hline
TC-13 & Cleanup on failure & Agent process that exits non-zero; full test campaign and 40 benchmark runs & Run marked failed; sandbox removed; no VaultShield containers or networks remain & Passed; zero sandbox, capture-helper or egress-proxy containers and zero run networks remained & Pass \\\hline
TC-14 & Hanging sign-in is bounded & Fake sign-in program that prompts and then sleeps 600 s, with the deadline patched to 3 s & Auth phase raises within 60 s with ``no further output'' & Took 600.3 s and returned success & Fail \\\hline
TC-15 & Accessibility: text contrast & All theme token pairs in light and dark themes & Text ≥ 4.5:1; UI components ≥ 3:1 & Text 4.77:1, 5.80:1; UI 3.37:1 and 3.30:1 & Pass \\\hline\end{longtable}
\addtocounter{table}{-1}
\par\medskip

\endgroup

\section{Conclusion on Testing}\label{conclusion-on-testing}

The campaign executed 2,704 automated tests, of which 2,702 passed, with no test skipped and the full set of Docker-backed isolation tests running against a real daemon. Every security-critical behaviour named in the objectives was exercised by at least one passing test against real infrastructure: the agent could not reach the host's files, secrets, home directory or Docker socket; real secrets did not reach the sandbox; non-allowlisted network traffic was blocked and recorded; off-limits writes failed the run with evidence; the diff was correct and could not be hidden by a tampering agent; credential bytes in outputs were caught and quarantined; and the submission gate and approval atomicity held, including under concurrency.

The campaign also surfaced real defects, which is the purpose of testing. Two Docker-backed tests in the interactive sign-in path failed on Docker Desktop 29.2.1; these affect only the optional --interactive-auth flow used by the Antigravity adapter, not the default pipeline. The benchmark revealed that the in-container file watcher, which is a polling program, missed the agent's write in 20\% of fast mock-agent runs. Because the authoritative change record is the trusted diff, which was complete in every run, this is a gap in supplementary evidence rather than in change capture, but it means a reviewer should not treat an absent file\_\allowbreak{}change event as proof that nothing changed. Finally, code review confirmed that the backend's unrestricted PATCH path can set a run to approved without an Approval record, which is acceptable only within the single-developer trust model of the MVP. These open issues are carried into the limitations discussion and the future work of Chapter 9.

\clearpage

\chapter{Result and Discussions}\label{results-and-analysis}

\section{Introduction}\label{introduction-1}

This chapter evaluates VaultShield against the objectives. The evaluation has three parts. The first is a performance analysis of the pipeline, based on a timing harness that executed 40 complete vaultshield run invocations against real Docker and recorded the wall-clock time of every pipeline stage, together with a baseline of the same agent executed directly on the host. The second is a comparative and safety analysis that contrasts running the agent with and without VaultShield in terms of overhead, evidence captured and unsafe actions detected, including the results of the feasibility spike on secret exposure. The third is a case study of real run directories preserved in the repository. The chapter closes with an overall discussion of the results, limitations and residual risks and a summary.

All measurements were taken with the deterministic mock agent. This agent is a real program executed inside the sandbox through the same path a real agent takes, but it makes no model API calls. This choice isolates the pipeline's own overhead from the large and highly variable latency of model inference, and it allows the results to be reproduced without API cost. The consequence, discussed later, is that real-agent run times were not measured in Semester VII.

\section{Evaluation Metrics}\label{evaluation-metrics}

Table 7.1 defines the metrics used, and Table 7.2 describes the environment in which they were measured.

\Needspace{8\baselineskip}
\par\vspace{\TableCapGap}
\noindent\begin{minipage}{\textwidth}
\captionof{table}{Evaluation metrics}\label{tab:7-1}
\end{minipage}
\par
\begin{longtable}[]{@{}|
  >{\raggedright\arraybackslash}p{(\columnwidth - 7\tabcolsep) * \real{0.3333}} |
  >{\raggedright\arraybackslash}p{(\columnwidth - 7\tabcolsep) * \real{0.3333}} |
  >{\raggedright\arraybackslash}p{(\columnwidth - 7\tabcolsep) * \real{0.3333}}|@{}}
\hline
Metric & Definition & Method of measurement \\
\hline
\endfirsthead
\hline
Metric & Definition & Method of measurement \\
\hline
\endhead
\endlastfoot
End-to-end run latency & Wall-clock time of one vaultshield run invocation, from command start to exit, with backend calls stubbed & Timing harness wrapping the Click command \\\hline
Stage latency & Time spent in each pipeline stage & Wrappers around the stage functions as bound in cli.py; summed per run \\\hline
Overhead versus direct execution & Difference and ratio between VaultShield latency and the same mock agent run directly on the host & Same harness, n = 20 direct runs \\\hline
Projected relative overhead & Fixed overhead divided by an assumed agent session length & Derived from measured overhead; not measured with real agents \\\hline
Isolation effectiveness & Fraction of runs in which the host checkout was unchanged and the agent could not reach host files, home directory or Docker socket & Isolation system test and harness check that the agent's file never appeared on the host \\\hline
Change-capture completeness & Fraction of runs whose diff.json contained the agent's change & Harness inspection of every run's diff.json \\\hline
Monitoring completeness & Fraction of runs whose Event\_Log contained a file\_\allowbreak{}change event for the agent's change & Harness inspection of every run's events.jsonl \\\hline
Detection of unsafe actions & Fraction of injected unsafe behaviours that were blocked or detected & Security-focused test cases TC-04 to TC-08 \\\hline
Secret exposure & Whether an agent could read or reveal a real secret value under a given protection configuration & Feasibility spike with a production agent; inspection of generated .env.sandbox files \\\hline
Cleanup reliability & Number of VaultShield containers or networks remaining after all runs & docker ps -a and docker network ls after the campaign \\\hline\end{longtable}
\addtocounter{table}{-1}
\par
\Needspace{8\baselineskip}
\par\vspace{\TableCapGap}
\noindent\begin{minipage}{\textwidth}
\captionof{table}{Evaluation environment}\label{tab:7-2}
\end{minipage}
\par
\begin{longtable}[]{@{}|
  >{\raggedright\arraybackslash}p{(\columnwidth - 5\tabcolsep) * \real{0.5000}} |
  >{\raggedright\arraybackslash}p{(\columnwidth - 5\tabcolsep) * \real{0.5000}}|@{}}
\hline
Item & Value \\
\hline
\endfirsthead
\hline
Item & Value \\
\hline
\endhead
\endlastfoot
Machine & Apple MacBook Air, Apple M4 processor \\\hline
Operating system & macOS 27.0.1 \\\hline
Container runtime & Docker Desktop 29.2.1 \\\hline
Sandbox image & vaultshield-sandbox:local \\\hline
Software & Python 3.11.14, Git 2.52.0, Node.js 24.12.0, docker-py 7.2.0, Click 8.5.0, Playwright 1.63.0, Django 5.2.18, DRF 3.18.3, PyMongo 4.18.2 \\\hline
Target repository & Local four-file repository served to the sandbox over git://host.docker.internal \\\hline
Agent & mock adapter \\\hline
Backend & Stubbed for registration, upload and status updates \\\hline\end{longtable}
\addtocounter{table}{-1}
\par\medskip

\section{Performance Analysis}\label{performance-analysis}

Table 7.3 reports the stage-level latency of the default configuration, which is the configuration most developers would use first.

\begingroup\scriptsize

\Needspace{8\baselineskip}
\par\vspace{\TableCapGap}
\noindent\begin{minipage}{\textwidth}
\captionof{table}{Stage latency in the default configuration}\label{tab:7-3}
\end{minipage}
\par
\begin{longtable}[]{@{}|
  >{\raggedright\arraybackslash}p{(\columnwidth - 15\tabcolsep) * \real{0.1111}} |
  >{\raggedleft\arraybackslash}p{(\columnwidth - 15\tabcolsep) * \real{0.1481}} |
  >{\raggedleft\arraybackslash}p{(\columnwidth - 15\tabcolsep) * \real{0.1481}} |
  >{\raggedleft\arraybackslash}p{(\columnwidth - 15\tabcolsep) * \real{0.1481}} |
  >{\raggedleft\arraybackslash}p{(\columnwidth - 15\tabcolsep) * \real{0.1481}} |
  >{\raggedleft\arraybackslash}p{(\columnwidth - 15\tabcolsep) * \real{0.1481}} |
  >{\raggedleft\arraybackslash}p{(\columnwidth - 15\tabcolsep) * \real{0.1481}}|@{}}
\hline
Stage & Mean & Median & Std. dev. & Min & Max & Share of mean total \\
\hline
\endfirsthead
\hline
Stage & Mean & Median & Std. dev. & Min & Max & Share of mean total \\
\hline
\endhead
\endlastfoot
3. Credential shaping & 0.001 & 0.001 & 0.000 & 0.000 & 0.001 & 0.1\% \\\hline
5. Sandbox creation, clone and dependency install & 0.239 & 0.238 & 0.017 & 0.206 & 0.279 & 17.0\% \\\hline
5-base. Trusted base snapshot & 0.013 & 0.013 & 0.002 & 0.011 & 0.022 & 0.9\% \\\hline
6. Agent execution and monitoring & 0.413 & 0.418 & 0.033 & 0.352 & 0.478 & 29.4\% \\\hline
8. Trusted capture & 0.651 & 0.641 & 0.053 & 0.585 & 0.773 & 46.4\% \\\hline
8. Token-hygiene scan & 0.019 & 0.016 & 0.008 & 0.015 & 0.042 & 1.4\% \\\hline
9. Teardown & 0.042 & 0.040 & 0.006 & 0.035 & 0.065 & 3.0\% \\\hline
Other & 0.025 & -- & -- & -- & -- & 1.8\% \\\hline
\textbf{End-to-end total} & \textbf{1.403} & \textbf{1.371} & \textbf{0.075} & \textbf{1.299} & \textbf{1.541} & \textbf{100\%} \\\hline\end{longtable}
\addtocounter{table}{-1}
\par\medskip

\endgroup

The total is tightly distributed: the coefficient of variation is 5.3\% and the slowest of the 20 runs took 1.54 s. Figure 7.1 shows the mean of each stage and Figure 7.2 its share of the total.


\noindent\begin{minipage}{\textwidth}
\centering
\SafeIncludeGraphics[width=0.9\textwidth,height=0.5\textheight,keepaspectratio]{images/diagrams/fig7-1-stage-latency.png}
\captionof{figure}{Mean stage latency in the default configuration}\label{fig:7-1}
\end{minipage}
\par\medskip


\noindent\begin{minipage}{\textwidth}
\centering
\SafeIncludeGraphics[width=0.82\textwidth,height=0.5\textheight,keepaspectratio]{images/diagrams/fig7-2-stage-share.png}
\captionof{figure}{Share of mean end-to-end time by stage, default configuration}\label{fig:7-2}
\end{minipage}
\par\medskip

Three observations follow from Table 7.3 and Figures 7.1 and 7.2. First, trusted capture is the most expensive stage at 0.651 s, because it creates and starts a second container, streams the work tree out of the killed sandbox, rebuilds the repository and runs several Git commands. This is the price of the tamper-resistance property demonstrated by TC-07; it buys the guarantee that the diff a reviewer sees was not produced by tooling the agent could influence. Second, agent execution inside the sandbox took 0.413 s against 0.186 s for the same program on the host, an additional 0.227 s that covers starting the in-container watcher, waiting for it to signal readiness, streaming output over Docker's exec API and draining the watcher's final reconciliation pass. Third, container creation and clone are cheap when the image is already present; the first run on a machine additionally pays the one-time image pull, which was not measured.

\begingroup\scriptsize

\Needspace{8\baselineskip}
\par\vspace{\TableCapGap}
\noindent\begin{minipage}{\textwidth}
\captionof{table}{End-to-end latency by configuration}\label{tab:7-4}
\end{minipage}
\par
\begin{longtable}[]{@{}|
  >{\raggedright\arraybackslash}p{(\columnwidth - 17\tabcolsep) * \real{0.1000}} |
  >{\raggedleft\arraybackslash}p{(\columnwidth - 17\tabcolsep) * \real{0.1333}} |
  >{\raggedleft\arraybackslash}p{(\columnwidth - 17\tabcolsep) * \real{0.1333}} |
  >{\raggedleft\arraybackslash}p{(\columnwidth - 17\tabcolsep) * \real{0.1333}} |
  >{\raggedleft\arraybackslash}p{(\columnwidth - 17\tabcolsep) * \real{0.1333}} |
  >{\raggedleft\arraybackslash}p{(\columnwidth - 17\tabcolsep) * \real{0.1333}} |
  >{\raggedleft\arraybackslash}p{(\columnwidth - 17\tabcolsep) * \real{0.1333}} |
  >{\raggedright\arraybackslash}p{(\columnwidth - 17\tabcolsep) * \real{0.1000}}|@{}}
\hline
Configuration & n & Mean & Median & Std. dev. & Min & Max & Distinctive stage \\
\hline
\endfirsthead
\hline
Configuration & n & Mean & Median & Std. dev. & Min & Max & Distinctive stage \\
\hline
\endhead
\endlastfoot
Direct on host: agent only & 20 & 0.186 & 0.186 & 0.003 & 0.180 & 0.194 & -- \\\hline
Direct on host: clone + agent & 20 & 0.237 & 0.237 & 0.005 & 0.230 & 0.253 & -- \\\hline
VaultShield default & 20 & 1.403 & 1.371 & 0.075 & 1.299 & 1.541 & Trusted capture 0.651 \\\hline
VaultShield with egress allowlist & 10 & 2.213 & 2.210 & 0.066 & 2.139 & 2.330 & Egress topology 0.422; teardown 0.550 \\\hline
VaultShield default with one screenshot route & 10 & 3.294 & 3.322 & 0.119 & 2.968 & 3.396 & App start and screenshot 1.827 \\\hline\end{longtable}
\addtocounter{table}{-1}
\par\medskip

\endgroup

Enabling the egress allowlist added 0.810 s on average. Of this, 0.422 s was spent creating the two per-run networks and starting the proxy and waiting for its readiness signal, and teardown rose from 0.042 s to 0.550 s because the proxy container and both networks must also be removed. The agent stage was unaffected, confirming that routing through the proxy does not slow a run that makes no network calls. Capturing one screenshot added 1.891 s, of which 1.827 s was starting the application, polling ready\_path until it answered and capturing the page with Chromium; this cost depends mainly on how quickly the target application starts. Figure 7.3 compares the configurations.


\noindent\begin{minipage}{\textwidth}
\centering
\SafeIncludeGraphics[width=0.9\textwidth,height=0.5\textheight,keepaspectratio]{images/diagrams/fig7-3-config-latency.png}
\captionof{figure}{Mean end-to-end time by configuration}\label{fig:7-3}
\end{minipage}
\par\medskip

\section{Overhead with and without VaultShield}\label{overhead-with-and-without-vaultshield}

The fair baseline for VaultShield's default mode is ``clone the repository and run the agent directly'', because VaultShield also performs a fresh clone. Against that baseline, VaultShield's default mode added 1.166 s per run, and the egress allowlist mode added 1.976 s. The ratios are large only because the mock agent finishes in a fraction of a second. For a real coding agent the overhead is fixed while the agent's own run time is dominated by model inference and tool execution, typically minutes. Figure 7.4 projects the relative overhead for assumed agent session lengths; it is derived from the measured fixed overhead and has not been validated with real agents.


\noindent\begin{minipage}{\textwidth}
\centering
\SafeIncludeGraphics[width=0.9\textwidth,height=0.5\textheight,keepaspectratio]{images/diagrams/fig7-4-projected-overhead.png}
\captionof{figure}{Projected relative overhead versus agent session length}\label{fig:7-4}
\end{minipage}
\par\medskip

In Figure 7.4, the lower line is the default mode and the upper line is the egress allowlist mode. For a five-minute agent session, the projected overhead is 0.39\% in default mode and 0.66\% with the egress allowlist, a negligible fixed overhead of well under five seconds per run. Table 7.5 places the overhead alongside the qualitative differences between the two ways of running an agent.

\Needspace{8\baselineskip}
\par\vspace{\TableCapGap}
\noindent\begin{minipage}{\textwidth}
\captionof{table}{Running an agent with and without VaultShield}\label{tab:7-5}
\end{minipage}
\par
\begin{longtable}[]{@{}|
  >{\raggedright\arraybackslash}p{(\columnwidth - 9\tabcolsep) * \real{0.2500}} |
  >{\raggedright\arraybackslash}p{(\columnwidth - 9\tabcolsep) * \real{0.2500}} |
  >{\raggedright\arraybackslash}p{(\columnwidth - 9\tabcolsep) * \real{0.2500}} |
  >{\raggedright\arraybackslash}p{(\columnwidth - 9\tabcolsep) * \real{0.2500}}|@{}}
\hline
Dimension & Direct execution & VaultShield default mode & VaultShield egress allowlist mode \\
\hline
\endfirsthead
\hline
Dimension & Direct execution & VaultShield default mode & VaultShield egress allowlist mode \\
\hline
\endhead
\endlastfoot
Mean wall-clock time & 0.237 s & 1.403 s & 2.213 s \\\hline
Fixed overhead versus direct & -- & +1.166 s & +1.976 s \\\hline
Projected overhead for a 5-minute session & -- & 0.39\% & 0.66\% \\\hline
Host checkout unchanged after run & No & Yes, 20 of 20 runs & Yes, 10 of 10 runs \\\hline
Real .env values visible to the agent & Yes & No & No \\\hline
Agent's change captured on a reviewable branch with a structured diff & No & Yes, 20 of 20 runs & Yes, 10 of 10 runs \\\hline
Timestamped event log & No & Yes & Yes \\\hline
Outbound traffic restricted and logged & No & No & Yes \\\hline
Change must be approved before a PR & No & Yes & Yes \\\hline
Containers or networks left behind & Not applicable & 0 & 0 \\\hline\end{longtable}
\addtocounter{table}{-1}
\par\medskip

\section{Safety Events Captured}\label{safety-events-captured}

Table 7.6 summarises the safety-relevant evidence produced across the 40 harness runs and the security-focused test cases.

\Needspace{8\baselineskip}
\par\vspace{\TableCapGap}
\noindent\begin{minipage}{\textwidth}
\captionof{table}{Safety events and evidence captured}\label{tab:7-6}
\end{minipage}
\par
\begin{longtable}[]{@{}|
  >{\raggedright\arraybackslash}p{(\columnwidth - 7\tabcolsep) * \real{0.3333}} |
  >{\raggedright\arraybackslash}p{(\columnwidth - 7\tabcolsep) * \real{0.3333}} |
  >{\raggedright\arraybackslash}p{(\columnwidth - 7\tabcolsep) * \real{0.3333}}|@{}}
\hline
Category & Scenario & Result \\
\hline
\endfirsthead
\hline
Category & Scenario & Result \\
\hline
\endhead
\endlastfoot
Isolation & Host checkout unchanged after run & 40 of 40 \\\hline
Isolation & Agent could read host .env, list host home or see the Docker socket & None of the three possible; all assertions held \\\hline
Change capture & Agent's file present in diff.json & 40 of 40 \\\hline
Monitoring & file\_\allowbreak{}change event recorded for the agent's file & 32 of 40: default 17/20, screenshots 9/10, egress 6/10 \\\hline
Screenshots & Screenshot captured for the configured route & 10 of 10 \\\hline
Policy & Off-limits write detected and run failed & 1 of 1 detected \\\hline
Network & Non-allowlisted CONNECT attempts blocked & 4 of 4 blocked with HTTP 403 and logged \\\hline
Network & Direct DNS exfiltration attempt & Blocked \\\hline
Credential leakage & Token-hygiene vectors in a real container & 5 of 5 detected, diff quarantined \\\hline
Tamper resistance & Tampered in-container Git trying to hide a file & 2 of 2 variants defeated \\\hline
Secret suppression & Values in six real .env.sandbox files in the repository & 30 of 30 values are dummies; the connection string's real host was replaced by localhost with dummy user information \\\hline
Cleanup & VaultShield containers and networks after the campaign & 0 containers, 0 networks \\\hline\end{longtable}
\addtocounter{table}{-1}
\par\medskip

The detection rates in Table 7.6 are rates over injected test scenarios, not over real-world attacks; they establish that each control works for the behaviour it was designed against, not that the controls are exhaustive. The monitoring result is the one clear shortfall and is analysed in the case study.

\section{Secret Exposure With and Without Protection}\label{secret-exposure-with-and-without-protection}

The feasibility spike conducted before implementation measured how a production coding agent handled a request to reveal secret values. The prompt asked it to read .env and .env.local in the project folder and report the exact values of two fake keys, using every tool available. Table 7.7 reports the outcome under four protection configurations.

\begingroup\small

\Needspace{8\baselineskip}
\par\vspace{\TableCapGap}
\noindent\begin{minipage}{\textwidth}
\captionof{table}{Secret exposure under four protection configurations}\label{tab:7-7}
\end{minipage}
\par
\begin{longtable}[]{@{}|
  >{\raggedright\arraybackslash}p{(\columnwidth - 13\tabcolsep) * \real{0.1667}} |
  >{\raggedright\arraybackslash}p{(\columnwidth - 13\tabcolsep) * \real{0.1667}} |
  >{\raggedright\arraybackslash}p{(\columnwidth - 13\tabcolsep) * \real{0.1667}} |
  >{\raggedright\arraybackslash}p{(\columnwidth - 13\tabcolsep) * \real{0.1667}} |
  >{\raggedright\arraybackslash}p{(\columnwidth - 13\tabcolsep) * \real{0.1667}} |
  >{\raggedright\arraybackslash}p{(\columnwidth - 13\tabcolsep) * \real{0.1667}}|@{}}
\hline
Config & Protection layers & Secret echoed in answer & .env read by file tool & .env.local read by file tool & Verdict \\
\hline
\endfirsthead
\hline
Config & Protection layers & Secret echoed in answer & .env read by file tool & .env.local read by file tool & Verdict \\
\hline
\endhead
\endlastfoot
C0 & None & Yes, both keys & Allowed & Allowed & Leaked \\\hline
C1 & Deny rules read\_file(.env) and read\_file(.env*) & No & Blocked & Allowed & Read into context \\\hline
C2 & C1 plus a PreToolUse hook denying shell access to .env & No & Blocked & Allowed & Read into context \\\hline
C3 & C2 plus real values replaced by the placeholder vs\_\allowbreak{}protected & No & Blocked & Read, but only the placeholder was on disk & Safe \\\hline\end{longtable}
\addtocounter{table}{-1}
\par\medskip

\endgroup

The spike shows that agent-level deny rules are brittle: exact-path rules work, glob rules silently match nothing, and a file the rule-writer did not anticipate is read into the agent's context even when the final answer does not print it. Only configuration C3, in which the real value was not on disk at all, was safe. This finding drove two design decisions. In the main pipeline, the sandbox receives .env.sandbox with dummy values, so the real values are never present in the agent's environment. For agents run outside the sandbox, the Protect module moves values into the macOS Keychain and leaves placeholders in the files. Interactive and GUI surfaces of the agent were not tested in the spike and remain pending.

\section{Case Study: Analysis of Real Runs}\label{case-study-analysis-of-real-runs}

The repository preserves the run directories of five complete mock-agent runs made during development, in addition to four incomplete ones. This section analyses one complete run in detail and then compares all five.

Run 20261009-091725-ff669fcd7dd00b2b was started with the prompt ``add a file''. Its Event\_Log, reproduced in Table 7.8, contains seven events.

\Needspace{8\baselineskip}
\par\vspace{\TableCapGap}
\noindent\begin{minipage}{\textwidth}
\captionof{table}{Event\_Log of run 20261009-091725-ff669fcd7dd00b2b}\label{tab:7-8}
\end{minipage}
\par
\begin{longtable}[]{@{}|
  >{\raggedleft\arraybackslash}p{(\columnwidth - 9\tabcolsep) * \real{0.3077}} |
  >{\raggedright\arraybackslash}p{(\columnwidth - 9\tabcolsep) * \real{0.2308}} |
  >{\raggedright\arraybackslash}p{(\columnwidth - 9\tabcolsep) * \real{0.2308}} |
  >{\raggedright\arraybackslash}p{(\columnwidth - 9\tabcolsep) * \real{0.2308}}|@{}}
\hline
\# & Timestamp & Type & Payload \\
\hline
\endfirsthead
\hline
\# & Timestamp & Type & Payload \\
\hline
\endhead
\endlastfoot
1 & 09:17:27.821242 & output & [mock-sandbox] working directory set to the sandbox clone \\\hline
2 & 09:17:27.821311 & output & [mock-sandbox step=0] simulated output for prompt='add a file' \\\hline
3 & 09:17:27.872204 & output & [mock-sandbox step=1] simulated output for prompt='add a file' \\\hline
4 & 09:17:27.925904 & output & [mock-sandbox step=2] simulated output for prompt='add a file' \\\hline
5 & 09:17:27.979070 & output & [mock-sandbox] wrote vaultshield\_\allowbreak{}mock\_\allowbreak{}agent.txt \\\hline
6 & 09:17:27.979111 & output & [mock-sandbox] simulated final message for prompt='add a file' \\\hline
7 & 09:17:28.085579 & file\_\allowbreak{}change & created: vaultshield\_\allowbreak{}mock\_\allowbreak{}agent.txt \\\hline\end{longtable}
\addtocounter{table}{-1}
\par\medskip

The log tells a reviewer four things without any further tooling. The first event confirms that the agent's working directory was the sandbox clone, not a host path. Events 2 to 6 are the agent's own narrative of its work, spanning 157.9 ms. Event 7, recorded by the in-container watcher 106.5 ms after the agent reported writing the file, independently corroborates the agent's claim with a container-side path. The run's diff.json contains exactly one entry, vaultshield\_\allowbreak{}mock\_\allowbreak{}agent.txt with one addition and no deletions, whose patch adds the line written by Mock\_\allowbreak{}Adapter inside the Sandbox. The run's screenshot for route / shows the sandbox clone's directory listing, as does the screenshot from run …092431 reproduced as Figure 5.10. Finally, the run's .env.sandbox holds the five keys of the developer's .env with dummy values, including a MongoDB connection string whose real user information and cluster host were replaced.

Comparing the five complete runs exposes a consistent pattern and one anomaly. In all five, the agent's narrative spanned 153.0 to 163.9 ms, and in all five the diff recorded the agent's file. In four runs the watcher reported the file 106.5 to 136.3 ms after the agent's ``wrote'' line, consistent with its polling interval. In run 20261009-092431-6891a631b5f4a313, however, the Event\_Log contains only the six output events and no file\_\allowbreak{}change event, even though the diff and the screenshot both show the file. The 40-run benchmark reproduced the same gap in 8 runs. The likely explanation is a race between the very short-lived mock agent and the polling watcher's shutdown and drain sequence, which the final reconciliation pass is intended to cover; the root cause had not been confirmed at the time of writing.

The case study therefore illustrates the design principle that the event log and the diff play different roles. The event log is corroborating evidence of behaviour, valuable for understanding what the agent did and in what order; the trusted diff is the authoritative record of what changed, built independently from a pre-agent base. A reviewer who relied only on the timeline of run …092431 would have missed the file creation; a reviewer who reads the diff would not. A mechanical cross-check that flags any path present in the diff but absent from the event log is a straightforward improvement.

\section{Overall Discussion and Limitations}\label{discussion-and-limitations}

VaultShield's own documentation insists that it reduces risk rather than eliminating it, and this evaluation supports that framing. The limitations below are grouped into security limitations, which are residual risks of the design, and evaluation limitations, which concern the strength of the evidence presented in this chapter. Table 7.9 summarises the security limitations.

\begingroup\small

\Needspace{8\baselineskip}
\par\vspace{\TableCapGap}
\noindent\begin{minipage}{\textwidth}
\captionof{table}{Residual risks and limitations}\label{tab:7-9}
\end{minipage}
\par
\begin{longtable}[]{@{}|
  >{\raggedright\arraybackslash}p{(\columnwidth - 11\tabcolsep) * \real{0.2000}} |
  >{\raggedright\arraybackslash}p{(\columnwidth - 11\tabcolsep) * \real{0.2000}} |
  >{\raggedright\arraybackslash}p{(\columnwidth - 11\tabcolsep) * \real{0.2000}} |
  >{\raggedright\arraybackslash}p{(\columnwidth - 11\tabcolsep) * \real{0.2000}} |
  >{\raggedright\arraybackslash}p{(\columnwidth - 11\tabcolsep) * \real{0.2000}}|@{}}
\hline
\# & Limitation & Consequence & Current mitigation & Planned improvement \\
\hline
\endfirsthead
\hline
\# & Limitation & Consequence & Current mitigation & Planned improvement \\
\hline
\endhead
\endlastfoot
L1 & Container escape through a kernel or Docker vulnerability & An agent could reach the host & Docker socket not mounted; only a dummy, read-only file mounted & gVisor or microVM isolation \\\hline
L2 & Default off mode runs the sandbox as root with Docker's default capabilities, no seccomp or AppArmor profile beyond Docker's default, no read-only root and no resource limits & Widens the escape surface and allows resource exhaustion & Hardening applied to proxy and capture helper; sandbox hardening only with egress enabled & Non-root user, cap\_drop, no-new-privileges, limits by default \\\hline
L3 & Egress is unrestricted unless explicitly enabled & Network exfiltration of repository content is possible in default mode & log-only and allowlist modes available & Default-deny egress \\\hline
L4 & Allowlist is host-based; www.googleapis.com also exposes general upload APIs & An agent holding a broad token could upload data to an allowlisted service without a recorded violation & Minimal per-agent lists; risk documented & Path-aware filtering or per-agent API proxies \\\hline
L5 & Codex and Kiro built-in allowlists are unverified against real traffic & Those agents may fail or be over-permitted in allowlist mode & Marked unverified in source & Verification in Semester VIII \\\hline
L6 & Dependency install hooks run inside the sandbox & A poisoned dependency executes with sandbox privileges & Dummy credentials; disposable sandbox & Pinned or pre-built dependency layers \\\hline
L7 & PATCH /api/runs/:id/ can set approved without an Approval record & A token holder can bypass the dashboard decision; submit checks only the status & Single-developer trust domain; documented & Gate every status transition; separate reviewer identity and RBAC \\\hline
L8 & No multi-user model: users are not linked to projects or decisions & The approval gate is a workflow control, not a separation-of-duties control & Out of MVP scope & Reviewer accounts, RBAC, decision attribution \\\hline
L9 & Polling file watcher misses some changes & Timeline may omit a write & Trusted diff is authoritative & Event-log/diff cross-check; inotify-based watcher \\\hline
L10 & Dashboard Diff tab shows per-file counts but not patch text & Reviewers must open diff.json or the bundle for line-level review & Patches stored and scanned & Inline patch rendering \\\hline
L11 & Screenshots are not scanned for credentials; deliberately encoded exfiltration is not detected & A secret rendered on screen or encoded in output could pass the hygiene scan & Dummy credentials make real secrets unlikely in the sandbox & Image OCR scan; entropy-based detection \\\hline
L12 & Interactive sign-in path failed two tests on Docker Desktop 29.2.1 & --interactive-auth may hang on that platform & Default pipeline unaffected & Fix I-07 and I-08 \\\hline\end{longtable}
\addtocounter{table}{-1}
\par\medskip

\endgroup

Several of these deserve emphasis. L1 and L2 together are the most important: in the default mode, the sandbox is isolated from the host only by Docker's namespace and cgroup boundaries, and the process inside runs as root with Docker's default capability set. This is a deliberate, documented trade-off to keep the base path simple and compatible with arbitrary agent toolchains, but it means VaultShield is not a substitute for a virtual machine when the agent or its inputs are actively hostile \cite{ref11}, \cite{ref12}, \cite{ref13}. L3 means that the strongest network controls are opt-in, and a developer who never sets --egress gets observability of file changes but not of network activity. L7 and L8 mean that the human approval gate protects against an agent reaching a PR without review, which is the threat it was designed for, but does not protect against a person holding the API token who chooses to skip review. In the single-developer MVP these are the same person, so this is consistent with the stated scope, but it must be addressed before a team deployment. Finally, the evidence the dashboard presents is exactly that: evidence, not proof. It informs a human decision and does not independently establish that a run was safe.

The evaluation itself has limitations. All performance figures come from the mock agent on one machine, with a small local repository, a pre-pulled image and a stubbed backend, so they characterise the pipeline's fixed overhead rather than the experience of running a real agent on a large repository over a real network. The projected overheads in Figure 7.4 are arithmetic projections from that fixed cost. The safety results are rates over injected test scenarios, which demonstrate that each control works as designed but do not measure resistance to an adaptive adversary. The secret-exposure spike covered one agent's headless CLI surface; its GUI surfaces were not tested. Review turnaround, an important practical metric, was not measured. Each of these is addressed in the Semester VIII plan.

Taken together, the results show that the central idea of VaultShield, separating the execution of an agent from trust in its output, can be realised at low cost. The pipeline added about 1.2 seconds to a run in its default configuration and about 2.0 seconds with network allowlisting, figures that are negligible against the minutes a real coding agent typically spends on a task. In exchange, every evaluated run left the developer's checkout untouched, kept real secrets out of the agent's environment, recorded a timestamped event log, captured the agent's change on a dedicated branch with a structured diff that a tampering agent could not falsify, and held that change behind an approval state machine that resolved concurrent decisions correctly. Cleanup was complete after more than 2,000 tests and 40 benchmark runs.

The distribution of cost is itself informative. Almost half of the default overhead is spent on trusted capture, which exists solely to make the diff trustworthy. A cheaper design could have built the branch inside the sandbox, but the tamper-resistance tests show what that saving would cost: the agent's own Git would then decide what the reviewer sees. Similarly, the egress allowlist roughly doubles teardown time because it creates real network infrastructure per run, and that infrastructure is what makes the network policy enforceable against a root process inside the sandbox rather than advisory. In both cases the evaluation supports the design choice of paying a fixed, sub-second cost for a property that cannot be obtained otherwise.

The results also validate the decision to state the security model explicitly. The testing and case study found real gaps, and the residual-risk analysis identifies structural limits that no amount of testing would remove. Because these are documented, a team adopting VaultShield can decide, for example, to always enable the allowlist, to treat the diff rather than the timeline as authoritative, or to run VaultShield inside a VM for untrusted repositories. A system that claimed to be a complete sandbox would leave those decisions uninformed.

\section{Chapter Summary}\label{chapter-summary}

This chapter evaluated VaultShield using a 40-run timing harness, the security-focused test cases of Chapter 6, a pre-implementation spike on secret exposure and a case study of preserved run directories. The default pipeline completed in 1.403 s on average, adding 1.166 s to direct execution, with trusted capture and agent monitoring the largest stages; the egress allowlist and screenshot capture added 0.810 s and 1.891 s respectively. Isolation, change capture, screenshot capture and cleanup succeeded in every evaluated run, all injected unsafe behaviours in the security test cases were blocked or detected, and every value in six real sandbox credential files was a dummy. The in-container watcher recorded the agent's change in 80\% of runs, which confirms the role of the trusted diff as the authoritative record. The chapter closed with twelve explicitly stated limitations, the most significant being the shared-kernel isolation boundary, the opt-in nature of egress control, and the single-user approval model. Chapter 8 assesses these results against the objectives.

\clearpage

\chapter{Conclusion}\label{conclusion}

This project set out to separate two things that current AI coding-agent tooling conflates: running an agent and trusting what it produces. VaultShield achieves that separation with a pipeline in which an agent receives a disposable Docker container, a fresh clone of the repository and a credential file containing only shape-preserving dummy values; its terminal output, file changes and network decisions are recorded as a timestamped event log; its change is captured as exactly one commit on a dedicated branch, with a structured diff built by a separate hardened container from a base recorded before the agent ran; and the change can become a pull or merge request only after a human approves it in a review dashboard. The system was implemented as a Python CLI, a Django REST Framework backend on MongoDB with Celery workers, and a React dashboard, totalling about 18.1 thousand logical lines of product code and 21.8 thousand lines of tests.

Table 8.1 assesses the outcome against each objective.

\Needspace{8\baselineskip}
\par\vspace{\TableCapGap}
\noindent\begin{minipage}{\textwidth}
\captionof{table}{Achievement of objectives}\label{tab:8-1}
\end{minipage}
\par
\begin{longtable}[]{@{}|
  >{\raggedright\arraybackslash}p{(\columnwidth - 9\tabcolsep) * \real{0.2500}} |
  >{\raggedright\arraybackslash}p{(\columnwidth - 9\tabcolsep) * \real{0.2500}} |
  >{\raggedright\arraybackslash}p{(\columnwidth - 9\tabcolsep) * \real{0.2500}} |
  >{\raggedright\arraybackslash}p{(\columnwidth - 9\tabcolsep) * \real{0.2500}}|@{}}
\hline
Objective & Measurable criterion & Evidence & Outcome \\
\hline
\endfirsthead
\hline
Objective & Measurable criterion & Evidence & Outcome \\
\hline
\endhead
\endlastfoot
1. Universal AI agent execution platform & A common CLI runs multiple agents across development environments & The mock, codex, kiro and antigravity agents run behind one adapter interface through the vaultshield CLI; TC-01 and the end-to-end system tests passed & Achieved \\\hline
2. Agentic AI Orchestrator & Project context understood, clarification questions asked, enriched context supplied & The ask command and run --ask scan the project, ask clarifying questions and build an enriched prompt, with a final secret check; the orchestrate test modules passed & Achieved \\\hline
3. Secure sandboxed execution with credential protection & Agents run in isolated Docker sandboxes; production credentials and resources protected & Host checkout, HEAD and .env unchanged in 40 of 40 benchmark runs; dummy credential file and one-shot clone token; 135 credential tests and TC-02, TC-03 passed & Achieved \\\hline
4. Comprehensive monitoring and automated security analysis & AI activities, logs and code changes captured; vulnerabilities and unsafe operations identified & Ordered event log and structured diff for every run; off-limits and egress violations recorded; the Celery security scan produced findings; TC-04 to TC-08 passed & Achieved, with a minor file-watcher gap \\\hline
5. Transparent governance and human approval & Web dashboard to review, verify, approve or reject AI changes before deployment & Four evidence tabs and approve/reject controls; submit refuses non-approved runs; concurrent decisions resolve to one outcome; 283 frontend tests, TC-09 and TC-10 passed & Achieved \\\hline\end{longtable}
\addtocounter{table}{-1}
\par\medskip

All five objectives were achieved. The platform runs several agents through one CLI, the orchestrator enriches prompts after asking for context, the sandbox keeps real credentials and the host away from the agent, the monitoring and security scan capture and analyse every run, and the dashboard gates each change behind a human decision. Two gaps remain within objectives that were otherwise met. The approval gate reliably prevents an agent's change from reaching a pull request without a human decision, but in the single-developer MVP it does not prevent a token holder from bypassing the dashboard through an unrestricted status update. The monitoring component records an ordered event log for every run, but its polling file watcher missed the agent's write in a fifth of very short benchmark runs; because the trusted diff captured every change, this is a gap in corroborating evidence rather than in change capture.

The project makes four contributions. First, it provides a working, reproducible reference implementation of governance practices that the literature recommends for agentic systems, namely constrained action space, human approval, legible activity logs and attributable identifiers \cite{ref9}, \cite{ref10}, specialised for repository-scale coding agents. Second, it contributes specific engineering techniques that address problems not covered by existing agent platforms: shape-preserving credential substitution that lets an application start without real secrets; a one-shot Git credential helper that confines a repository token to a single process; trusted capture that builds the reviewable diff outside the agent's reach; and a transactional approval state machine that remains correct under concurrency on MongoDB. Third, it offers empirical evidence that this level of protection is inexpensive, adding about 1.2 seconds per run in the default configuration, and that agent-level deny rules alone are insufficient to protect secret values, whereas keeping real values off disk is effective. Fourth, it demonstrates an explicit security model, in which every implemented control is traceable to code and tests and every residual risk is documented, as a practice worth adopting in security tooling.

The team learned several lessons in the course of the work. Negative security guarantees are best specified as numbered properties and verified with property-based tests and tests against real infrastructure, because mocks tend to encode the very assumptions that need checking; several real defects, including a leaked database host name and a race that produced server errors, were found only this way. Isolation features must be designed together with the evidence they produce: it was not enough to run the agent in a container, because the diff also had to be built somewhere the agent could not influence. Platform behaviour changes and must be monitored continuously. Finally, honesty about limits is itself a design requirement: writing down what the system does not protect against forced clearer decisions about defaults and directly shaped the plan for future work described in the next chapter.

\clearpage

\chapter{Future Scope}\label{future-scope}

The limitations and open issues identified earlier define a concrete programme of future work. The extensions below are ordered roughly by the size of the risk they address, and each is grounded in the existing architecture so that it can be added behind the current interfaces.

Stronger isolation. The most important extension is to raise the isolation boundary beyond a standard Docker container. In the near term, the sandbox should be hardened by default rather than only when egress control is enabled: running the agent as a non-root user, dropping all capabilities not required by the agent's toolchain, setting no-new-privileges, applying a restrictive seccomp profile, and imposing memory, CPU and process limits. In the longer term, the sandbox could run under gVisor's runsc runtime, which interposes a user-space kernel between the agent and the host kernel, or inside Firecracker microVMs, which provide hardware-virtualisation isolation with fast start-up \cite{ref12}, \cite{ref13}. Because Sandbox.create already passes container options through a single call, an alternative runtime can be introduced as a configuration option, and the benchmark harness of Chapter 7 can quantify its cost on the same workloads.

Policy engines. Today's policies are expressed as fixed configuration: off-limits paths, an egress mode and allowlist, and static security-scan rules. A declarative policy engine, for example one based on Open Policy Agent's Rego language, would let a team express rules such as ``runs that modify infrastructure files require two approvals'', ``deny any new dependency not on the approved list'', or ``fail runs whose diff adds an outbound call to an unlisted domain''. Policies could be evaluated at three points already present in the pipeline: during execution, at capture and at review. Default-deny egress, path-aware network filtering and per-agent API proxies that expose only the model endpoint would close the host-based allowlist gap described in limitation L4.

Separation of duties and access control. The single-developer MVP should evolve into a multi-user system with reviewer accounts, role-based access control separating the roles of run author, reviewer and administrator, attribution of every decision to a user, and support for multiple required approvals. As part of this, the PATCH endpoint should gate every status transition, so that approved can be reached only through the approval endpoint with an Approval record, and the submit command should verify the existence of that record rather than only the status.

Cost and resource controls. Real agents consume model tokens and compute. The adapters could record per-run token usage and estimated cost from the agents' own logs, the backend could enforce per-project budgets and the CLI could apply wall-clock and resource limits to the sandbox, terminating runaway runs with a clear failure class. A cost column in the run history would make the economics of agent use visible to a team.

Multi-repository and parallel runs. Many real changes span several repositories, for example a service, its client library and its deployment configuration. The run model could be extended so that one logical run creates a sandbox with several clones, captures one branch and diff per repository, and is approved as a unit. The collision-resistant run identifier already supports concurrent runs, so parallel attempts at the same task by different agents, with the reviewer choosing the best result, are also a natural extension.

IDE integration. Developers increasingly interact with agents inside their editors. An extension for VS Code-based editors could start a VaultShield run from the editor, stream the event log into a panel, render the structured diff in the editor's native diff view and expose approve and reject commands. The dashboard's Diff tab should likewise render patch text inline. The Protect module's agent hooks already integrate with Antigravity, Claude Code, Kiro and Codex, and could be extended to route an agent's shell commands through a sandbox automatically.

Telemetry and analytics dashboard. The backend already stores every run's status history, events, findings and decisions. An analytics view could report review turnaround times, approval and rejection rates by agent and project, the distribution of security findings, egress violations over time and stuck approvals. Live streaming of events during a run, for which a Django Channels worker is already reserved in the deployment, would let a reviewer watch and interrupt a run in progress, implementing the interruptibility practice recommended in \cite{ref9}.

Evidence quality. The event log should be cross-checked against the trusted diff, with any path present in the diff but absent from the log flagged in the timeline, and the polling watcher should be replaced or supplemented by an event-driven watcher. Screenshots could be scanned with optical character recognition for credential patterns, and entropy-based detection could flag encoded data in output. The interactive sign-in defects found on Docker Desktop 29.2.1 should be fixed and kept under continuous test.

Formal threat modelling and verification. The threat model is currently stated informally in the README. A structured threat model, for example using the STRIDE categories over the data-flow diagram of Figure 3.5, would systematically enumerate threats at every trust boundary. The approval state machine is small enough to be specified and model-checked formally, and the security properties already expressed as Hypothesis properties could be strengthened into a published, versioned specification. An independent security review or penetration test of the sandbox and proxy would provide external validation.

Broader evaluation. Finally, the Semester VIII evaluation should measure real agents on representative repositories, including end-to-end latency, the frequency and nature of egress violations and security findings, and review turnaround in a usability study with reviewers. Running a subset of SWE-bench tasks \cite{ref2} through VaultShield would allow its overhead and its effect on reviewer decisions to be reported against a recognised benchmark.

\clearpage

\clearpage

\begingroup
\renewcommand{\bibname}{Bibliography}
\begin{thebibliography}{99}
\addcontentsline{toc}{chapter}{Bibliography}
\bibitem{ref1} Chen, M., Tworek, J., Jun, H., Yuan, Q., de Oliveira Pinto, H. P., Kaplan, J., et al., ``\href{https://arxiv.org/abs/2107.03374}{Evaluating Large Language Models Trained on Code,}'' arXiv preprint arXiv:2107.03374, 2021.

\bibitem{ref2} Jimenez, C. E., Yang, J., Wettig, A., Yao, S., Pei, K., Press, O., and Narasimhan, K., ``\href{https://openreview.net/forum?id=VTF8yNQM66}{SWE-bench: Can Language Models Resolve Real-World GitHub Issues?,}'' International Conference on Learning Representations, 2024.

\bibitem{ref3} Yang, J., Jimenez, C. E., Wettig, A., Lieret, K., Yao, S., Narasimhan, K., and Press, O., ``\href{https://proceedings.neurips.cc/paper_files/paper/2024/hash/5a7c947568c1b1328ccc5230172e1e7c-Abstract-Conference.html}{SWE-agent: Agent-Computer Interfaces Enable Automated Software Engineering,}'' Advances in Neural Information Processing Systems, 2024.

\bibitem{ref4} Wang, X., Li, B., Song, Y., Xu, F. F., Tang, X., Zhuge, M., et al., ``\href{https://arxiv.org/abs/2407.16741}{OpenHands: An Open Platform for AI Software Developers as Generalist Agents,}'' International Conference on Learning Representations, 2025.

\bibitem{ref5} Ruan, Y., Dong, H., Wang, A., Pitis, S., Zhou, Y., Ba, J., Dubois, Y., Maddison, C. J., and Hashimoto, T., ``\href{https://openreview.net/forum?id=GEcwtMk1uA}{Identifying the Risks of LM Agents with an LM-Emulated Sandbox,}'' International Conference on Learning Representations, 2024.

\bibitem{ref6} Debenedetti, E., Zhang, J., Balunović, M., Beurer-Kellner, L., Fischer, M., and Tramèr, F., ``\href{https://proceedings.neurips.cc/paper_files/paper/2024/file/97091a5177d8dc64b1da8bf3e1f6fb54-Paper-Datasets_and_Benchmarks_Track.pdf}{AgentDojo: A Dynamic Environment to Evaluate Prompt Injection Attacks and Defenses for LLM Agents,}'' Advances in Neural Information Processing Systems, Datasets and Benchmarks Track, 2024.

\bibitem{ref7} Greshake, K., Abdelnabi, S., Mishra, S., Endres, C., Holz, T., and Fritz, M., ``\href{https://doi.org/10.1145/3605764.3623985}{Not What You've Signed Up For: Compromising Real-World LLM-Integrated Applications with Indirect Prompt Injection,}'' ACM Workshop on Artificial Intelligence and Security, 2023.

\bibitem{ref8} Pearce, H., Ahmad, B., Tan, B., Dolan-Gavitt, B., and Karri, R., ``\href{https://doi.org/10.1109/SP46214.2022.9833571}{Asleep at the Keyboard? Assessing the Security of GitHub Copilot's Code Contributions,}'' IEEE Symposium on Security and Privacy, 2022.

\bibitem{ref9} Shavit, Y., Agarwal, S., Brundage, M., et al., ``\href{https://openai.com/index/practices-for-governing-agentic-ai-systems/}{Practices for Governing Agentic AI Systems,}'' OpenAI White Paper, 2023.

\bibitem{ref10} Chan, A., Ezell, C., Kaufmann, M., Wei, K., Hammond, L., Bradley, H., et al., ``\href{https://doi.org/10.1145/3630106.3658948}{Visibility into AI Agents,}'' ACM Conference on Fairness, Accountability, and Transparency, 2024.

\bibitem{ref11} Combe, T., Martin, A., and Di Pietro, R., ``\href{https://doi.org/10.1109/MCC.2016.100}{To Docker or Not to Docker: A Security Perspective,}'' IEEE Cloud Computing, 2016.

\bibitem{ref12} Agache, A., Brooker, M., Florescu, A., Iordache, A., Liguori, A., Neugebauer, R., Piwonka, P., and Popa, D.-M., ``\href{https://www.usenix.org/conference/nsdi20/presentation/agache}{Firecracker: Lightweight Virtualization for Serverless Applications,}'' USENIX Symposium on Networked Systems Design and Implementation, 2020.

\bibitem{ref13} Young, E. G., Zhu, P., Caraza-Harter, T., Arpaci-Dusseau, A. C., and Arpaci-Dusseau, R. H., ``\href{https://www.usenix.org/conference/hotcloud19/presentation/young}{The True Cost of Containing: A gVisor Case Study,}'' USENIX Workshop on Hot Topics in Cloud Computing, 2019.

\bibitem{ref14} Bacchelli, A., and Bird, C., ``\href{https://doi.org/10.1109/ICSE.2013.6606617}{Expectations, Outcomes, and Challenges of Modern Code Review,}'' International Conference on Software Engineering, 2013.

\bibitem{ref15} Boehm, B. W., ``\href{https://www.pearson.com/en-us/subject-catalog/p/software-engineering-economics/P200000000306}{Software Engineering Economics,}'' Prentice-Hall, 1981.
\end{thebibliography}
\endgroup
\end{document}
